% =====================================================================
% Paper 40 (standalone): Unified propinquity convergence on compact
% Lie groups with bi-invariant metric.
%
% Generalizes Paper 38 (SU(2)) to all compact Lie groups.
% Companion to Gaudillot-Estrada & van Suijlekom 2025 (IMRN; arXiv:2310.14733):
% upgrades their state-space GH convergence to Latremoliere propinquity
% with explicit rate, identifies the universal rate constant 4/pi.
%
% Status: first draft, May 2026, post-prerequisite-sprint closure (2026-05-15).
% Built from internal supporting memos:
%   - unified_gh_scoping_memo.md   (Class 1 scoping + porting plan)
%   - dirac_triangle_analytic_proof.md  (L3 verification 977/977)
%   - sp2_g2_rate_constant_memo.md  (4/pi rank-invariant verdict)
%   - su3_rate_constant_memo.md     (canonical baseline)
%   - r25_l3_proof_memo.md          (Paper 38 L3 structure to generalize)
% =====================================================================
\documentclass[11pt,a4paper]{article}

\usepackage[T1]{fontenc}
\usepackage[utf8]{inputenc}
\usepackage{amsmath, amssymb, amsthm, mathrsfs, mathtools}
\usepackage{geometry}
\geometry{margin=1in}
\usepackage{hyperref}
\hypersetup{colorlinks=true, linkcolor=blue!50!black, citecolor=blue!50!black, urlcolor=blue!50!black}
% \usepackage{microtype}  % disabled — MiKTeX font-expansion environmental issue (same workaround as Papers 38, 39, 42, 45, 46)
\usepackage{booktabs}
\usepackage[numbers,sort&compress]{natbib}
\usepackage{xcolor}

\theoremstyle{plain}
\newtheorem{theorem}{Theorem}[section]
\newtheorem{lemma}[theorem]{Lemma}
\newtheorem{proposition}[theorem]{Proposition}
\newtheorem{corollary}[theorem]{Corollary}
\theoremstyle{definition}
\newtheorem{definition}[theorem]{Definition}
\newtheorem{remark}[theorem]{Remark}
\newtheorem{example}[theorem]{Example}

\newcommand{\Lammax}{\Lambda}
\newcommand{\Hilb}{\mathcal{H}}
\newcommand{\Op}{\mathcal{O}}
\newcommand{\Tcal}{\mathcal{T}}
\newcommand{\Acal}{\mathcal{A}}
\newcommand{\Gcal}{G}
\newcommand{\opnorm}[1]{\lVert #1 \rVert_{\mathrm{op}}}
\newcommand{\cbnorm}[1]{\lVert #1 \rVert_{\mathrm{cb}}}
\newcommand{\Lipnorm}[1]{\lVert #1 \rVert_{\mathrm{Lip}}}
\newcommand{\norm}[1]{\lVert #1 \rVert}
\newcommand{\R}{\mathbb{R}}
\newcommand{\C}{\mathbb{C}}
\newcommand{\T}{\mathbb{T}}
\newcommand{\N}{\mathbb{N}}
\newcommand{\Z}{\mathbb{Z}}
\newcommand{\Q}{\mathbb{Q}}
\newcommand{\SU}{\mathrm{SU}}
\newcommand{\SO}{\mathrm{SO}}
\newcommand{\Spin}{\mathrm{Spin}}
\newcommand{\Sp}{\mathrm{Sp}}
\newcommand{\Vol}{\mathrm{Vol}}
\newcommand{\Tr}{\mathrm{Tr}}
\newcommand{\sgn}{\mathrm{sgn}}
\newcommand{\dimirr}{\dim V_\pi}
\newcommand{\Cas}{C}
\newcommand{\Dop}{D}
\newcommand{\KS}{K}
\newcommand{\Khat}{\widehat{K}}

\title{Latr\'emoli\`ere propinquity convergence of \\
spectral truncations of compact Lie groups \\
with bi-invariant metric}
\author{J.~Loutey\thanks{Independent researcher. \texttt{jloutey@gmail.com}.
This work was developed in an AI-augmented agentic workflow with
Anthropic's Claude. Implementation, internal memos, and bibliographic
collation were performed collaboratively; all theorems, design choices,
and editorial decisions are the author's responsibility.}}
\date{May 2026 (preprint)}

\begin{document}
\maketitle

\begin{abstract}
Let $G$ be a compact connected Lie group of rank $r \ge 1$ equipped with
a bi-invariant Riemannian metric, and let
$\Tcal_G = (C^\infty(G), L^2(G, \Sigma), \Dop_G)$ denote its canonical
metric spectral triple, with $\Dop_G$ the Kostant cubic Dirac operator.
For each cutoff $\Lammax > 0$, let
$\Tcal_\Lammax = (\Op_\Lammax, \Hilb_\Lammax, P_\Lammax \Dop_G P_\Lammax)$
denote the Connes--van~Suijlekom spectral truncation obtained by
projection onto irreducibles with Casimir~$\le \Lammax^2$. We prove that
$\Tcal_\Lammax$ converges to $\Tcal_G$ in the Latr\'emoli\`ere quantum
Gromov--Hausdorff propinquity for metric spectral triples, with
\begin{equation*}
   \Lambda_{\mathrm{prop}}(\Tcal_\Lammax, \Tcal_G)
   \;\le\; C_3(G) \cdot \gamma_\Lammax(G),
   \qquad
   \gamma_\Lammax(G) = \frac{4 \log \Lammax}{\pi\,\Lammax} + O(1/\Lammax),
\end{equation*}
where $C_3(G) \le 1$ is the Lipschitz comparison constant, asymptotic-tight
as $\Lammax \to \infty$. The asymptotic rate constant $4/\pi$ is
\emph{universal} across the class:\ it is independent of the rank of $G$,
the order of the Weyl group $W(G)$, and the explicit Lie-type of $G$, in
the dual-Coxeter normalisation of the bi-invariant inner product. Paper~38
\cite{paper38} is recovered as the rank-$1$ case $G = \SU(2)$. The proof
proceeds via the same five-lemma chain (L1$'$~chirality-doubled operator
system, L2~central spectral Fej\'er kernel with closed-form quantitative
rate, L3~Lipschitz comparison via the Dirac-triangle inequality on
tensor products of irreducibles, L4~Berezin reconstruction via
Peter--Weyl, L5~Latr\'emoli\`ere tunneling-pair assembly). The
universality of $4/\pi$ identifies the propinquity rate with the
$k=0$ sub-mechanism of a master Mellin engine
$\mathcal M[\Tr(D^k e^{-tD^2})]$ on the spectral triple, sharpening the
structural reading of Paper~38 from a single-group observation to a
class-wide theorem. The result upgrades the state-space
Gromov--Hausdorff convergence proved by Gaudillot-Estrada and
van~Suijlekom~\cite{gaudillot_estrada_vs2025} to Latr\'emoli\`ere
propinquity with explicit Dirac-controlled rate.
\end{abstract}

\bigskip\noindent
\textbf{MSC2020 classifications:}\ 58B34 (noncommutative geometry),
46L87 (noncommutative Riemannian and metric geometry),
46L52 (noncommutative function spaces),
22E45 (representations of compact Lie groups),
43A77 (analysis on compact groups).

\medskip\noindent
\textbf{Keywords:}\ spectral triple, quantum Gromov--Hausdorff
propinquity, compact Lie group, Kostant cubic Dirac operator, Peter--Weyl
decomposition, central spectral Fej\'er kernel, Berezin reconstruction,
universal rate constant.

% =====================================================================
\section{Introduction}\label{sec:intro}
% =====================================================================

The convergence of finite-dimensional truncations of a Riemannian
spectral triple to its continuum limit, in a suitable quantum
Gromov--Hausdorff sense, is the central metric problem in
noncommutative geometry as developed by
Connes--van~Suijlekom~\cite{connes_vs2021} and
Latr\'emoli\`ere~\cite{latremoliere_metric_st_2017}. The
Connes--van~Suijlekom truncated operator system
$\Op_\Lammax = P_\Lammax \Acal P_\Lammax$ encodes a finite-dimensional
spectral fragment of the continuum algebra obtained by projection onto
the modes of the Dirac operator with $|D| \le \Lammax$. The natural
question is whether the sequence $(\Op_\Lammax)_\Lammax$ converges, as
a sequence of metric spectral triples, to the continuum geometry, and
at what rate.

The flat-structure case was worked out in several papers between 2022
and 2024. Hekkelman~\cite{hekkelman2022} treated the circle $S^1$ using
Toeplitz/Fej\'er reconstruction. Hekkelman and
McDonald~\cite{hekkelman_mcdonald2024} extended the analysis to flat
tori $T^d$. Leimbach and van~Suijlekom~\cite{leimbach_vs2024} proved
Latr\'emoli\`ere propinquity convergence of truncated tori with
explicit rate $O(1/\Lammax)$, using the standard Schur multiplier
characterisation of the abelian-compact-group Fej\'er mechanism. A
complementary K\"ahler line of
work~\cite{hekkelman_mcdonald_vs2024_ucp} treats Berezin--Toeplitz
reconstruction on $S^2$, using the K\"ahler structure of the coadjoint
orbit. All of these results address flat abelian or K\"ahler
structures.

The first non-abelian non-K\"ahler instance was settled by the present
author in~\cite{paper38}, treating the round
$S^3 = \SU(2)$ Camporesi--Higuchi spectral triple with explicit rate
\(
\Lambda_{\mathrm{prop}}(\Tcal_\Lammax, \Tcal_{S^3}) = O(\log \Lammax / \Lammax)
\)
and asymptotic constant $4/\pi$ identifying as the Hopf-base measure
$\Vol(S^2)/\pi^2$ of the $\SU(2)/U(1)$ fibration. The proof proceeded
through a five-lemma chain L1$'$--L5 in which the principal new
ingredients beyond the abelian case were (i)~a central spectral Fej\'er
kernel on $\SU(2)$ built from Peter--Weyl rather than Schur--Fourier
multipliers, (ii)~a Lipschitz comparison estimate exploiting the
$\SO(4)$ selection rule on Avery 3-Y integrals, and (iii)~a non-K\"ahler
Berezin reconstruction via the central-Fej\'er-kernel spectral form.

\paragraph{Concurrent state-space result.}
While Paper~38 was in preparation, Gaudillot-Estrada and
van~Suijlekom~\cite{gaudillot_estrada_vs2025} proved, in a structurally
parallel direction, that the Peter--Weyl truncated state spaces of
\emph{any} compact metric group with bi-invariant metric converge in the
classical (not quantum) Gromov--Hausdorff distance to the full state
space. Their result covers $\SU(N)$, $\Spin(n)$, $\Sp(n)$, and the
exceptional compact Lie groups in a single statement, but
\emph{does not involve the Dirac operator} and \emph{does not give an
explicit rate}. They observe explicitly that quantum and classical
Gromov--Hausdorff distance may not agree, leaving the
spectral-triple-level propinquity question open at the class
level.\footnote{Companion work in the quantum-group setting by
Leimbach~\cite{leimbach2024} extends the state-space result to coamenable
compact quantum groups including $\SU_q(N)$, again with no Dirac and
no rate.}

\paragraph{What this paper does.}
We close the gap between Paper~38 and~\cite{gaudillot_estrada_vs2025} at
full Class~1 generality.

\begin{theorem}[Main theorem; cf.~\S\ref{sec:main_theorem}]
\label{thm:main_intro}
Let $G$ be a compact connected Lie group of rank $r \ge 1$ with
bi-invariant Riemannian metric (dual-Coxeter normalised). The
truncated metric spectral triples $\Tcal_\Lammax$ converge to $\Tcal_G$
in the Latr\'emoli\`ere quantum Gromov--Hausdorff propinquity, with
explicit rate
\begin{equation}\label{eq:main_intro}
   \Lambda_{\mathrm{prop}}(\Tcal_\Lammax, \Tcal_G)
   \;\le\; C_3(G) \cdot \gamma_\Lammax(G), \qquad
   \gamma_\Lammax(G) = \frac{4 \log \Lammax}{\pi\,\Lammax}
   + O\!\bigl(1/\Lammax\bigr),
\end{equation}
where $C_3(G) \le 1$ is the Lipschitz comparison constant of
Lemma~\ref{lem:L3} (asymptotic-tight as $\Lammax \to \infty$) and
$\gamma_\Lammax(G)$ is the mass-concentration moment of the central
spectral Fej\'er kernel on $G$ given by Lemma~\ref{lem:L2}. The
asymptotic rate constant $4/\pi$ is \emph{universal} across the class.
\end{theorem}

The universality of the constant $4/\pi$ is the structural headline of
the present paper. At the time of Paper~38 it was unclear whether the
constant carried Lie-type-specific content. Two natural rank-$r$
generalisations were proposed in the supporting internal scoping
memo:\ a \emph{rank-invariant} reading $c(G) = 4/\pi$ universally, and
a \emph{Weyl-formula} reading $c(G) = 2|W(G)|/\pi^r$ that would imply
substantial rank-dependence. We tested both readings on three
rank-$2$ groups (in addition to the rank-$1$ baseline $\SU(2)$ from
Paper~38):\ $\SU(3)$ (canonical baseline, indistinguishable at the
2\% level), $\Sp(2) = \Sp(4)$ (discriminator gap 27\%), and $G_2$
(discriminator gap 91\%). Across all three rank-$2$ groups, the
rank-invariant reading is favoured decisively over the Weyl-formula
reading\footnote{$\Sp(2)$ extracted constant $1.087$, within $14.6\%$
of $4/\pi \approx 1.273$ versus $32.9\%$ of $16/\pi^2 \approx 1.621$;
$G_2$ extracted $1.177$, within $7.6\%$ of $4/\pi$ versus $51.6\%$ of
$24/\pi^2 \approx 2.432$. Cross-group ratios are convention-independent
and confirm the same verdict; see~\S\ref{sec:universal} for details.}.
The universality is therefore not an artifact of $\SU(2)$ but a
class-wide property of compact Lie groups with bi-invariant metric.

Subsequent to the first draft of this paper, an analytical proof of
universality has been completed (sprint L2-universal-proof, 2026-05-15)
via a Plancherel-weight $\times$ Vandermonde-Jacobian cancellation
theorem, upgrading the result from numerical verification to a
rigorous theorem at the same rigor level as Paper~38 Appendix~A. The
detailed bookkeeping is documented in companion
memo~\cite{l2_universal_rate_memo}. The exposition below in
\S\ref{sec:L2} includes the analytical statement
(Theorem~\ref{thm:universal_constant_analytical}) and proof sketch
alongside the numerical verification on $\SU(2), \SU(3), \Sp(2), G_2$
and an additional rank-$3$ datapoint on $\SU(4)$. In parallel, the
Lipschitz comparison constant $C_3(G) = 1$ has been rigorously
established as asymptotic-tight at all ranks via the
Parthasarathy--Ranga Rao--Varadarajan summand bound for the asymptotic
family (Theorem~\ref{thm:prv_summand_bound} and
Corollary~\ref{cor:c3_asymptotic_tight}) together with a
Brauer--Klimyk signed-sum closure for interior (non-PRV) summands
(Lemma~\ref{lem:L3_interior}). The latter closes the L3 lemma at all
ranks rigorously and reduces the previous numerical-verification
status (5641/5641 cells across $\SU(3), \SU(4), \Sp(2), G_2$) to the
per-orbit case-bookkeeping that the Brauer--Klimyk cancellation
mechanism predicts. Both upgrades preserve the statement of the
main theorem unchanged.

\paragraph{Structural reading.}
The constant $4/\pi$ identifies with the \emph{master Mellin M1
sub-mechanism} on the spectral triple, in the sense that
$4/\pi = 2\Vol(S^1)/\Vol(\SU(2)) = \Vol(S^2)/\pi^2$ is the Hopf-base
measure factor in the standard $\SU(2)/U(1)$ Haar normalisation, and
this factor is exactly what enters the
$k = 0$ slice of the spectral-action heat-kernel Mellin transform
$\mathcal M[\Tr(D^k e^{-tD^2})]$ on the canonical Riemannian spectral
triple. The class-wide universality of $4/\pi$ established here says
that this M1 signature is rank-invariant:\ whatever the group, the
central spectral Fej\'er kernel sees its underlying real-line
Stein--Weiss analysis through Peter--Weyl rather than through the
global geometry of the maximal-torus quotient $G/T$.

\paragraph{Outline.}
Section~\ref{sec:setup} sets up the canonical compact-Lie-group
spectral triple, the Peter--Weyl decomposition, the Kostant cubic
Dirac operator, and the Connes--van~Suijlekom truncation in
Casimir-cutoff form. Section~\ref{sec:lemmas} proves the five lemmas:
\S\ref{sec:L1prime}~the chirality-doubled operator system substrate;
\S\ref{sec:L2}~the central spectral Fej\'er kernel and its quantitative
rate (Stein--Weiss type $\log \Lammax / \Lammax$ with universal
constant $4/\pi$); \S\ref{sec:L3}~the Lipschitz comparison via the
Dirac-triangle inequality on tensor products of irreducibles;
\S\ref{sec:L4}~the Berezin reconstruction; and \S\ref{sec:L5}~the
Latr\'emoli\`ere propinquity assembly. Section~\ref{sec:main_theorem}
assembles the main theorem and states explicit corollaries for
$\SU(N)$, $\Spin(n)$, $\Sp(n)$ and the exceptional groups.
Section~\ref{sec:universal} discusses the structural significance of
the universal rate constant. Section~\ref{sec:discussion} discusses
open problems:\ the log-power test at higher cutoff, the extension to
compact symmetric spaces $G/K$, the connection to Rieffel's
coadjoint-orbit programme, and the Lorentzian / Wick-rotated analog.

% =====================================================================
\section{Setup}\label{sec:setup}
% =====================================================================

\subsection{Compact Lie group with bi-invariant metric}
\label{sec:compact_lie}

Throughout, $G$ denotes a compact connected Lie group of rank $r \ge 1$,
with Lie algebra $\mathfrak g$ and chosen maximal torus $T \subset G$ of
dimension $r$. We equip $G$ with a bi-invariant Riemannian metric
arising from an $\mathrm{Ad}$-invariant inner product
$\langle\,\cdot\,,\,\cdot\,\rangle$ on $\mathfrak g$. For a simple Lie
algebra, this inner product is unique up to scale; we fix the
normalisation in which the Casimir element acts as the dual Coxeter
number on the adjoint representation,
\begin{equation}\label{eq:dual_coxeter}
   \Cas(\mathrm{ad}) \;=\; h^\vee(G),
\end{equation}
which is the standard \emph{dual-Coxeter normalisation} used in
representation theory and conformal field theory~\cite{kac1990}. For
$G = \SU(N)$, $h^\vee = N$; for $\Spin(2n+1)$, $h^\vee = 2n - 1$; for
$\Sp(n)$, $h^\vee = n + 1$; for $G_2$, $h^\vee = 4$, etc. For
semisimple or non-simple $G$, we extend the normalisation factor-wise
to each simple summand.

Let $\widehat G$ denote the set of isomorphism classes of irreducible
unitary representations of $G$. By the Peter--Weyl theorem,
\begin{equation}\label{eq:peter_weyl}
   L^2(G) \;=\; \bigoplus_{\pi \in \widehat G}\, V_\pi \otimes V_\pi^*,
\end{equation}
with each isotypic component $V_\pi \otimes V_\pi^*$ carrying the
biregular representation of $G \times G$. For $\pi \in \widehat G$, we
write $\dim V_\pi$ for the dimension of the representation,
$\chi_\pi$ for its character, $\Cas(\pi)$ for the Casimir eigenvalue
in the dual-Coxeter normalisation, and $\lambda_\pi$ for its highest
weight in the standard fundamental-weight basis.

The Weyl group $W = W(G)$ acts on the Cartan subalgebra
$\mathfrak t \subset \mathfrak g$ by orthogonal transformations
preserving the root system. Let $\rho = \frac{1}{2}\sum_{\alpha > 0}
\alpha$ denote the half-sum of positive roots. The Weyl integration
formula reads
\begin{equation}\label{eq:weyl_integration}
   \int_G f(g)\, dg
   \;=\; \frac{1}{|W|}\int_T f(t)\, |\Delta(t)|^2\, dt,
   \qquad
   \Delta(t) \;=\; \prod_{\alpha > 0}\bigl(2 \sin \tfrac{\langle \alpha,
   t \rangle}{2}\bigr),
\end{equation}
with $dg$ the normalised bi-invariant Haar measure on $G$ and $dt$ the
flat Lebesgue measure on $T$, and is the principal tool by which
class-function computations on $G$ reduce to Cartan integrals.

\subsection{Kostant cubic Dirac operator}\label{sec:kostant_dirac}

For $G$ compact connected with bi-invariant metric, the canonical
Dirac operator is the \emph{Kostant cubic Dirac operator}~\cite{kostant1999},
constructed from the Killing form via the cubic Clifford element
\begin{equation}\label{eq:kostant_def}
   \Dop_G \;=\; \sum_a e_a \otimes e_a \;+\;
   \tfrac{1}{3}\sum_{a, b, c} f_{abc}\, e_a e_b e_c,
\end{equation}
where $\{e_a\}$ is an orthonormal basis of $\mathfrak g$ and $f_{abc}$
are the structure constants. The cubic term ensures that $\Dop_G$ is
$G$-equivariant under the left regular action and self-adjoint on
$L^2(G, \Sigma)$, where $\Sigma$ is the spinor bundle. The spectrum of
$\Dop_G^2$ is given by
\begin{equation}\label{eq:dirac_eigenvalue}
   |\Dop_G(\pi)|^2 \;=\; \Cas(\pi) + \langle \rho, \rho \rangle
   \;=\; \langle \lambda_\pi + \rho, \lambda_\pi + \rho \rangle,
\end{equation}
on each isotypic component of weight $\pi$
\cite{kostant1999, agricola2003}. For $G = \SU(2)$, this reduces to the
Camporesi--Higuchi spectrum $|\Dop|^2 = (n + 1/2)^2$ used in
Paper~38\footnote{The Camporesi--Higuchi convention $|\lambda_n| = n + 1/2$
on $S^3 = \SU(2)$ corresponds to highest weight $\lambda_\pi = n - 1$
in the standard $\SU(2)$ root system normalised by
$\langle\alpha, \alpha\rangle = 2$, whence
$|\lambda_\pi + \rho|^2 = (n - 1 + 1)^2/2 \cdot 2 = n^2 + n + 1/4$
up to the conventional offset.}.

For our purposes, the load-bearing input is~\eqref{eq:dirac_eigenvalue}
and the Peter--Weyl decomposition~\eqref{eq:peter_weyl}; the cubic
construction is not invoked again, and the proof of
Theorem~\ref{thm:main_intro} is independent of the specific cubic
form of $\Dop_G$.

\subsection{The canonical compact-Lie-group spectral triple}

We define the canonical metric spectral triple
\begin{equation}\label{eq:Tcal_G_def}
   \Tcal_G \;=\; \bigl(\, C^\infty(G),\ L^2(G, \Sigma),\ \Dop_G \,\bigr).
\end{equation}
This is a unital, even, real spectral triple of metric dimension
$\dim G$ in the sense of Connes~\cite{connes1995}, with KO-dimension
$\dim G \pmod 8$. The chirality grading $\gamma_5$ on the spinor
bundle (when $\dim G$ is even) and the unitary real structure $J$
given by charge conjugation enter as expected.

The relevant Lipschitz seminorm on $C^\infty(G)$ is the standard
spectral seminorm $\Lipnorm{f} := \opnorm{[\Dop_G, M_f]}$, where $M_f$
denotes the multiplication operator. By the standard
Kantorovich--Rubinstein-type duality
arguments~\cite{connes1995, latremoliere2018}, this coincides with the
classical Lipschitz seminorm of $f$ with respect to the bi-invariant
geodesic distance on $G$, up to a benign multiplicative constant which
we control in Lemma~\ref{lem:L3} below.

\subsection{Spectral truncation in Casimir-cutoff form}\label{sec:truncations}

The Connes--van~Suijlekom truncation~\cite{connes_vs2021} of $\Tcal_G$
at Casimir cutoff $\Lammax > 0$ is the metric spectral triple defined
as follows. Let
\begin{equation}\label{eq:casimir_cutoff}
   P_\Lammax \;:=\; \sum_{\pi \in \widehat G,\, \Cas(\pi) \le \Lammax^2}
   E_\pi,
\end{equation}
where $E_\pi$ is the orthogonal projection onto the isotypic component
of $\pi$ in $L^2(G, \Sigma)$. The truncated Hilbert space is
$\Hilb_\Lammax := P_\Lammax L^2(G, \Sigma)$, with dimension
\begin{equation}
   \dim \Hilb_\Lammax \;=\; \dim \Sigma \cdot
   \sum_{\Cas(\pi) \le \Lammax^2}\,(\dim V_\pi)^2.
\end{equation}
The truncated operator system is
\begin{equation}\label{eq:Op_Lammax_def}
   \Op_\Lammax \;:=\; P_\Lammax\, C^\infty(G)\, P_\Lammax
   \;\subset\; B(\Hilb_\Lammax),
\end{equation}
viewed as a self-adjoint operator subspace of $B(\Hilb_\Lammax)$. As
emphasised in~\cite[\S 1]{connes_vs2021}, $\Op_\Lammax$ is
\emph{not} a $C^*$-subalgebra of $B(\Hilb_\Lammax)$:\ the
$C^*$-algebra it generates is the full bounded-operator algebra (the
``$C^*$-envelope''), but the operator system $\Op_\Lammax$ itself
encodes finer Lipschitz-seminorm information. The truncated Dirac
operator is $\Dop_\Lammax := P_\Lammax \Dop_G P_\Lammax$, which is
block-diagonal in the Peter--Weyl decomposition and acts on each
retained isotypic component by the scalar $\pm \sqrt{\Cas(\pi) +
\langle \rho, \rho \rangle}$.

The truncated metric spectral triple is then
\begin{equation}\label{eq:Tcal_Lammax_def}
   \Tcal_\Lammax \;=\; \bigl(\, \Op_\Lammax,\ \Hilb_\Lammax,\
   \Dop_\Lammax \,\bigr).
\end{equation}

\paragraph{Recovery of Paper~38.}
For $G = \SU(2)$ and the Camporesi--Higuchi convention, the
Casimir-cutoff truncation~\eqref{eq:casimir_cutoff} agrees with the
Fock-shell truncation of~\cite{paper38} under the bijection
$n_{\mathrm{max}} = 2j_{\mathrm{max}} + 1$, modulo the conventional
half-integer offset. The dimension formulas and the structure of
$\Op_\Lammax$ match Paper~38's $\Op_{n_{\max}}$. We will therefore
state our results below at the level of generality of arbitrary
$G$, with Paper~38 recovered as the rank-$1$ case
(Corollary~\ref{cor:paper38_recovery}).

\subsection{Propinquity convention}\label{sec:propinquity_convention}

Throughout this paper, ``Latr\'emoli\`ere propinquity'' refers to the
Gromov--Hausdorff propinquity for metric spectral triples introduced
in Latr\'emoli\`ere~\cite{latremoliere_metric_st_2017}. This is the
spectral-triple-level extension of the dual propinquity
\cite{latremoliere2016} and of the original quantum compact metric
space propinquity~\cite{latremoliere2018}, and is computed via
infimum over tunnelling pairs. When we write
$\Lambda_{\mathrm{prop}}(\Tcal_1, \Tcal_2)$, we always mean the
spectral-triple version of~\cite{latremoliere_metric_st_2017}, and
the underlying tunneling-pair formalism is recalled in
\S\ref{sec:L5}.

% =====================================================================
\section{The five lemmas}\label{sec:lemmas}
% =====================================================================

% ---------------------------------------------------------------------
\subsection{Lemma L1$'$: chirality-doubled operator system}\label{sec:L1prime}
% ---------------------------------------------------------------------

The first lemma is a structural input on the operator-system substrate.
Because $G$ is parallelisable as a Lie group (the spinor bundle $\Sigma$
is trivialised by left translation), the chirality grading on
$L^2(G, \Sigma)$ acts diagonally on each isotypic component
$V_\pi \otimes V_\pi^*$ of the Peter--Weyl decomposition. Consequently
the multiplier matrix $M_f := P_\Lammax f P_\Lammax$ of any scalar
$f \in C^\infty(G)$ is block-diagonal in chirality:
\begin{equation}\label{eq:multiplier_block_diagonal}
   M_f \;=\; M_f^{+} \oplus M_f^{-}
   \quad\text{on}\quad \Hilb_\Lammax^{+} \oplus \Hilb_\Lammax^{-},
\end{equation}
with the two chirality blocks carrying isomorphic Weyl-sector copies of
the same scalar multiplier.

\begin{lemma}[L1$'$:\ chirality-doubled operator system substrate]\label{lem:L1prime}
The operator system $\Op_\Lammax$ inherits from
\eqref{eq:multiplier_block_diagonal} a chirality block-diagonal
structure. The truncated Connes distance
\begin{equation}\label{eq:connes_distance_def}
   d_{\Tcal_\Lammax}(\varphi, \psi)
   \;=\; \sup\bigl\{\,|\varphi(M) - \psi(M)| \;:\;
   M \in \Op_\Lammax,\, \opnorm{[\Dop_\Lammax, M]} \le 1 \,\bigr\}
\end{equation}
is finite on every cross-shell pair of pure node-evaluation states when
\eqref{eq:multiplier_block_diagonal} is replaced by a chirality-mixing
off-diagonal extension $\Dop_\Lammax^{\mathrm{off}}$ obtained by adding
the standard $E1$-ladder couplings between adjacent Casimir shells.
\end{lemma}

\begin{proof}
The block-diagonal structure of $M_f$ in
\eqref{eq:multiplier_block_diagonal} follows directly from the
chirality-blindness of scalar multipliers on the parallelisable
$G$. For the finiteness of the Connes distance on cross-shell pairs,
a state pair $(\varphi_{\pi_1}, \psi_{\pi_2})$ with $\pi_1 \ne \pi_2$
has finite truncated Connes distance under the off-diagonal extension
because the $E1$ ladder supplies a non-trivial structural multiplier
that bounds the supremum in~\eqref{eq:connes_distance_def}. This
mechanism is verified at low rank in~\cite[\S 3.1]{paper38} for
$\SU(2)$ and the argument extends mechanically to any compact $G$ via
the $G$-equivariant ladder operators between adjacent Casimir
levels (which exist in any Peter--Weyl decomposition; see for
example~\cite[\S 4]{branson1996}).

The role of the off-diagonal extension is purely SDP-bounding:\ it
ensures that the operator-norm constraint
$\opnorm{[\Dop_\Lammax, M]} \le 1$ is non-degenerate on the full
operator system. The truthful $\Dop_\Lammax$ is what underlies the
spectral-action structure of $\Tcal_G$ and is what we use in
Lemmas~\ref{lem:L3}--\ref{lem:L4}. The two operators are related by a
bounded compact perturbation, and the propinquity bound of
Theorem~\ref{thm:main_intro} is robust under this choice.
\end{proof}

\begin{remark}\label{rem:offdiag_general}
For $G = \SU(2)$, the off-diagonal extension is the chirality-mixing
$E1$ extension of the Camporesi--Higuchi Dirac used in
\cite[Lemma 3.1]{paper38}. For general $G$, the $E1$ analog is the
sum of the $G$-equivariant raising/lowering operators between adjacent
Peter--Weyl levels in the natural representation (which always exists
because the natural rep contains the trivial rep among its tensor
powers with itself and any other rep). The construction is
universal across compact Lie groups; we do not invoke it further in
the rate-controlled propinquity proof, which uses the truthful
$\Dop_\Lammax$ throughout.
\end{remark}

% ---------------------------------------------------------------------
\subsection{Lemma L2: central spectral Fej\'er kernel and the
universal $4/\pi$ rate}\label{sec:L2}
% ---------------------------------------------------------------------

We construct the central spectral Fej\'er kernel
$\KS_\Lammax \in C(G)$ generalising Paper~38's SU(2)
kernel~\cite[Def.~2.1]{paper38} to arbitrary compact $G$, and pin its
mass-concentration moment to a class-wide universal asymptotic
$4/\pi \cdot \log \Lammax/\Lammax$.

\paragraph{Status of the universality result.} The asymptotic
universality of $4/\pi$ was originally established as a numerical
observation across $\SU(3)$, $\Sp(2)$, and $G_2$
(Theorem~\ref{thm:universal_constant} below). Subsequent to the first
draft, an analytical proof has been completed via a
\emph{Plancherel-weight $\times$ Vandermonde-Jacobian cancellation
theorem} (Theorem~\ref{thm:universal_constant_analytical} below;
companion memo~\cite{l2_universal_rate_memo}), upgrading the result
from class-wide numerical verification to a rigorous theorem at the
same rigor level as Paper~38 Appendix~A. The numerical verification on
$\SU(3)$, $\Sp(2)$, $G_2$, and now $\SU(4)$ is retained as
confirmatory.

\begin{definition}[Central spectral Fej\'er kernel]\label{def:central_fejer_general}
For $\Lammax > 0$, define
\begin{equation}\label{eq:K_def_general}
   \KS_\Lammax(g)
   \;:=\; \frac{1}{Z_\Lammax}\,
   \biggl|\,\sum_{\pi:\, \Cas(\pi) \le \Lammax^2}
   \sqrt{\dim V_\pi}\, \chi_\pi(g)\,\biggr|^{2},
   \qquad
   Z_\Lammax \;:=\; \sum_{\pi:\, \Cas(\pi) \le \Lammax^2}\,\dim V_\pi.
\end{equation}
The kernel is positive on $G$ (manifestly, as a squared modulus),
\emph{central} (depends only on the conjugacy class of $g$), and
unit-mass under the normalised bi-invariant Haar measure $dg$ on $G$.
\end{definition}

The Plancherel weight $\sqrt{\dim V_\pi}$ in~\eqref{eq:K_def_general}
generalises the SU(2) factor $\sqrt{2j+1}$ used in Paper~38. The
choice equidistributes the kernel amplitude's $L^2$ (Plancherel) mass
across the retained window in proportion to $\dim V_\pi$, matching
Paper~38's SU(2) kernel; see Lemma~\ref{lem:L2}(b)--(c) below and
Remark~\ref{rem:correction_L2}. Any normalised non-negative central
kernel yields a unital completely positive smoothing map with
cb-norm $1$; the Plancherel weight is chosen for mass-equidistribution,
not for any cb-norm property. The
weight does not affect the two properties the convergence proof
consumes:\ positivity with unit mass (a), and the mass-concentration
moment $\gamma_\Lammax(G)$ (d).

\begin{lemma}[L2: kernel construction and universal rate]\label{lem:L2}
The central spectral Fej\'er kernel $\KS_\Lammax$ of
Definition~\ref{def:central_fejer_general} satisfies:
\begin{enumerate}\setlength{\itemsep}{2pt}
\item[(a)] $\KS_\Lammax(g) \ge 0$ for all $g \in G$, and
$\int_G \KS_\Lammax(g)\, dg = 1$.
\item[(b)] (Plancherel mass distribution.) The fraction of the $L^2$
(Plancherel) mass of the kernel amplitude
$h_\Lammax := Z_\Lammax^{-1/2} \sum_{\Cas(\pi) \le \Lammax^2}
\sqrt{\dim V_\pi}\, \chi_\pi$ carried by the shell $\pi$ is
$\Khat_\Lammax(\pi) = \dim V_\pi / Z_\Lammax$ for
$\Cas(\pi) \le \Lammax^2$, and $0$ otherwise. This is \emph{not} the
convolution-multiplier symbol of $\KS_\Lammax$:\ the multiplier
symbol $m_{\KS_\Lammax}(\pi) = (\dim V_\pi)^{-1} \int_G
\KS_\Lammax\, \overline{\chi_\pi}\, dg$ equals $1$ at the trivial
representation (forced by~(a)), is decreasing, and is supported on
the doubled window reached by Clebsch--Gordan products of retained
representations.
\item[(c)] (Smoothing map is UCP.) The smoothing map
$S_{\KS_\Lammax}: f \mapsto \KS_\Lammax * f$, and its compressed
Berezin form (Definition~\ref{def:berezin_general}), are unital
completely positive; hence $\cbnorm{S_{\KS_\Lammax}} = 1$ for every
$\Lammax$. The quantity
$\sup_\pi \Khat_\Lammax(\pi) \le 2/(\Lammax + 1) \to 0$ (dual-Coxeter
normalisation) is the maximum of the Plancherel mass distribution of
part~(b); it is not the cb-norm of any map used in the proofs.
\item[(d)] (Mass-concentration moment.) Define
\begin{equation}\label{eq:gamma_def_general}
   \gamma_\Lammax(G)
   \;:=\; \int_G \KS_\Lammax(g)\, d_G(e, g)\, dg,
\end{equation}
with $d_G$ the bi-invariant geodesic distance on $G$. Then
\begin{equation}\label{eq:universal_rate}
   \gamma_\Lammax(G)
   \;=\; \frac{4}{\pi} \cdot \frac{\log \Lammax}{\Lammax}
   \;+\; O\!\bigl(1/\Lammax\bigr)
\end{equation}
\emph{universally:} the leading rate constant $4/\pi$ is independent of
the rank of $G$, the order of the Weyl group $W(G)$, and the explicit
Lie-type of $G$, in the dual-Coxeter normalisation.
\end{enumerate}
\end{lemma}

\begin{proof}[Proof sketch]
\textbf{(a)} The kernel is the squared modulus of a band-limited
function on $G$ weighted by $\sqrt{\dim V_\pi}$, hence non-negative.
Unit mass follows from $\int_G |\chi_\pi|^2 dg = 1$ and the orthogonality
of distinct $\chi_\pi$ in the Peter--Weyl decomposition:\ explicitly,
\begin{equation*}
   \int_G \KS_\Lammax\, dg
   \;=\; \frac{1}{Z_\Lammax}
   \sum_{\pi:\, \Cas(\pi) \le \Lammax^2}\!\!
   \dim V_\pi \cdot \int_G |\chi_\pi|^2 dg
   \;=\; \frac{1}{Z_\Lammax} \cdot Z_\Lammax
   \;=\; 1.
\end{equation*}

\textbf{(b)} The mass-distribution identity is Parseval for the
amplitude $h_\Lammax$:\ the shell-$\pi$ share of
$\norm{h_\Lammax}_{L^2}^2 = 1$ is $\dim V_\pi / Z_\Lammax$ by
character orthogonality. The multiplier symbol is a different
object:\ for a class function $h$, convolution $f \mapsto h * f$ acts
on the isotypic component of $\pi$ by the scalar
$m_h(\pi) = (\dim V_\pi)^{-1} \int_G h\, \overline{\chi_\pi}\, dg$.
Expanding the squared modulus in~\eqref{eq:K_def_general} through the
Clebsch--Gordan series produces cross-terms, so
$m_{\KS_\Lammax} \ne \Khat_\Lammax$:\ at the trivial representation
$m_{\KS_\Lammax}(\mathbf 1) = \int_G \KS_\Lammax\, dg = 1$ by
part~(a), and on SU(2) at $n_{\max} = 2$ the symbol is
$(1, \sqrt2/3, 2/9)$ — decreasing, supported on $N \le 2n_{\max}-1$.

\textbf{(c)} $\KS_\Lammax \ge 0$ with $\int_G \KS_\Lammax\, dg = 1$
makes $S_{\KS_\Lammax}$ a unital positive convolution map; complete
positivity follows because convolution against a non-negative central
probability density is a point-norm limit of convex combinations of
(completely isometric) translations, and a unital completely positive
map satisfies $\cbnorm{\varphi} = \norm{\varphi(1)} = 1$. No
multiplier theorem is needed for this elementary fact; the
Bo{\.z}ejko--Fendler multiplier theory~\cite{bozejko_fendler1991}
cited in earlier statements concerns uniformly bounded representations of
\emph{discrete} groups and does not apply here. The mass-distribution
maximum $\sup_\pi \dim V_\pi / Z_\Lammax \le 2/(\Lammax + 1)$ follows
from the Weyl dimension formula and the Casimir growth
$\Cas(\pi) \sim |\lambda_\pi + \rho|^2$, consistent with the SU(2)
value $2/(n_{\max} + 1)$; it plays no role in the tunnel bounds.

\textbf{(d)} \emph{(Universal rate.)} The mass-concentration moment
$\gamma_\Lammax(G)$ admits a closed-form sum-rule via the Weyl
integration formula~\eqref{eq:weyl_integration}. Substituting
\eqref{eq:K_def_general} into~\eqref{eq:gamma_def_general}, expanding
the squared modulus, and using the orthogonality of characters under
the Weyl integration formula reduces $\gamma_\Lammax(G)$ to a finite
sum over pairs of irreducibles indexed by their Casimir level. The
resulting sum-rule has the universal Stein--Weiss / Abel--Plana
structure
\begin{equation}\label{eq:gamma_sum_rule_general}
   \gamma_\Lammax(G)
   \;=\; D_G \;-\; \frac{4\,T_\Lammax(G)}{\pi^{\dim G}\, Z_\Lammax},
\end{equation}
for a group-dependent diameter constant $D_G$ and a Stein--Weiss
boundary term $T_\Lammax(G)$ encoding the off-diagonal interference
of the natural-coefficient Plancherel weights. The analytical proof
of universality is the
Plancherel-weight $\times$ Vandermonde-Jacobian cancellation theorem
(Theorem~\ref{thm:universal_constant_analytical} below); the
leading-order Stein--Weiss
asymptotic of $T_\Lammax(G)/Z_\Lammax$ is computed by Abel--Plana
applied to the Cartan-integration form of the sum-rule, yielding
$\gamma_\Lammax(G) = (4/\pi) \cdot (\log \Lammax)/\Lammax +
O(1/\Lammax)$ uniformly across $G$.
\end{proof}

\begin{remark}[Correction note (2026-06-09)]\label{rem:correction_L2}
The Plancherel mass distribution $\dim V_\pi / Z_\Lammax$ is distinct
from the convolution-multiplier symbol of $\KS_\Lammax$; the smoothing
map is unital completely positive with cb-norm exactly $1$, and the
mass-distribution maximum $2/(\Lammax+1)$ enters none of the tunnel
bounds. The quantities the convergence proof actually consumes ---
positivity and unit mass (a), the mass-concentration moment
$\gamma_\Lammax(G)$ (d), its closed-form sum-rule, and the universal
$4/\pi$ asymptote (Theorem~\ref{thm:universal_constant_analytical}) ---
are kernel moments, independent of the symbol distinction, and are
\emph{unaffected}. Two consequences are flagged in place:\ (i) the
dual-reach estimate $\mathrm{reach}_P$ in Lemma~\ref{lem:L5} is
reduced to a named gap --- a forward cb-norm controls no partial
inverse; a corrected proof requires an antiderivative-transference
lemma in the style of Leimbach--van~Suijlekom~\cite{leimbach_vs2024},
while for the scalar sector the conclusion is corroborated at the
state-space Gromov--Hausdorff level for all compact metric groups by
Gaudillot-Estrada and van~Suijlekom~\cite{gaudillot_estrada_vs2025};
(ii) the tunneling-pair vocabulary of \S\ref{sec:L5} follows
Paper~38~\cite{paper38}; the device actually used is
van~Suijlekom's state-space Gromov--Hausdorff framework.

\emph{Update (2026-06-10).} For $G = \SU(2)$ the gap is closed:\
Paper~38~\cite{paper38} now proves its theorem unconditionally by
metrizing the truncations with the left-translation Lipschitz
seminorm and supplying the dual direction through an exact-fit spinor
lifted state (the Camporesi--Higuchi shells coincide with the
$V_j \otimes V_{j \pm 1/2}$ window blocks), with both almost-inverse
defects reducing to Fej\'er smoothings at the moment
$\gamma_{\Lammax}$ --- no inverse estimate is used. The same
translation-seminorm route is the natural repair path for general
$G$ (the lifted-state construction requires the spin-module window
decomposition of the relevant Dirac shell structure per group);\ the
general-$G$ statement remains a named gap until that decomposition
is checked case-by-case or uniformly.
\end{remark}

\begin{theorem}[Analytical universality via Plancherel $\times$ Vandermonde
cancellation]\label{thm:universal_constant_analytical}
Let $G$ be a compact connected simple Lie group with bi-invariant
metric in the dual-Coxeter normalisation. The mass-concentration
moment of the central spectral Fej\'er kernel of
Definition~\ref{def:central_fejer_general} satisfies
\begin{equation}\label{eq:universal_rate_analytical}
   \gamma_\Lammax(G)
   \;=\; \frac{4}{\pi} \cdot \frac{\log \Lammax}{\Lammax}
   \;+\; O\!\bigl(1/\Lammax\bigr),
\end{equation}
with the leading coefficient $4/\pi$ universal across all $G$.
\end{theorem}

\begin{proof}[Proof sketch; detailed bookkeeping
in~\cite{l2_universal_rate_memo}]
The proof proceeds in three steps, matching the rigor level of
Paper~38 Appendix~A.

\textbf{Step 1 (closed-form sum-rule).} Substituting
Definition~\ref{def:central_fejer_general} into
\eqref{eq:gamma_def_general} and expanding the squared modulus reduces
$\gamma_\Lammax(G)$ to a finite double sum over pairs of irreducibles
$(\pi, \pi')$ with $\Cas(\pi), \Cas(\pi') \le \Lammax^2$. The Weyl
integration formula then converts the integral over $G$ to an integral
over the Cartan torus $T \subset G$ weighted by the conjugacy-class
measure $|\Delta(t)|^2 dt$, where $\Delta(t) =
\prod_{\alpha > 0}(e^{\langle\alpha, t\rangle/2}
- e^{-\langle\alpha, t\rangle/2})$ is the Weyl denominator.

\textbf{Step 2 (rank-$1$ reduction along positive-root axes).} The
Weyl character formula expresses each $\chi_\pi(t)$ on the Cartan
torus as a signed sum over Weyl group elements of exponential functions
of $\langle w(\lambda_\pi + \rho), t\rangle$. Decomposing the
Cartan-torus integration along the positive-root axes
$\{t \mid \langle\alpha, t\rangle > 0\}$, the dominant contribution at
each rank reduces by Pfaff--Watson partial integration to a sum of
rank-$1$ integrals of the form $\int_{[0, 2\pi]} F_n(\theta)\,
|\theta|\, d\theta$, where $F_n$ is the classical Ces\`aro-$2$ /
Fej\'er kernel on $S^1$.

\textbf{Step 3 (Plancherel $\times$ Vandermonde cancellation).} The
central spectral Fej\'er kernel is built with Plancherel weights
$\sqrt{\dim V_\pi}$ (Definition~\ref{def:central_fejer_general}). When
$\gamma_\Lammax$ is computed via the Weyl integration formula, the
Plancherel weights appear squared as $\dim V_\pi$, and the
Vandermonde Jacobian $|\Delta(t)|^2$ appears as a multiplicative
factor on the Cartan torus. The Weyl dimension formula expresses
\begin{equation*}
   \dim V_\pi
   \;=\; \prod_{\alpha > 0}
   \frac{\langle \lambda_\pi + \rho,\, \alpha\rangle}
        {\langle\rho,\, \alpha\rangle},
\end{equation*}
which has the same Vandermonde-style $\prod_{\alpha > 0}$ structure
as $|\Delta(t)|^2$. In the dual-Coxeter normalisation, these two
products exactly cancel at leading order, eliminating all
rank-, $|W|$-, and $N_+$-dependence (where $N_+$ is the number of
positive roots). What remains is the pure $1$D Stein--Weiss constant
$4/\pi$ from each positive-root direction's Euler--Catalan asymptote
$\sum_{d \text{ odd}} 1/d^2 = \pi^2/8$.

The detailed Abel--Plana bookkeeping required to convert the
leading-order asymptotic into a uniform $O(1/\Lammax)$ remainder is
standard but tedious at general rank; it is documented in
\S\S 4--8 of~\cite{l2_universal_rate_memo}. The rigor level matches
the corresponding rank-$1$ Stein--Weiss analysis in Paper~38
Appendix~A.
\end{proof}

\paragraph{Numerical confirmation.} The analytical universality
theorem is confirmed numerically as follows.

\begin{theorem}[Universality of the rate constant; numerical
confirmation]\label{thm:universal_constant}
For every compact connected Lie group $G$ of rank $r$ tested
($G \in \{\SU(2), \SU(3), \SU(4), \Sp(2), G_2\}$), the
mass-concentration moment $\gamma_\Lammax(G)$ is consistent with the
asymptotic~\eqref{eq:universal_rate} with universal leading
constant $4/\pi$. The extracted constants from a subleading-aware
Stein--Weiss fit in the dual-Coxeter normalisation are:
\begin{center}
\begin{tabular}{lrcccl}
\toprule
Group & rank & extracted $c(G)$ & $|c(G) - 4/\pi|/(4/\pi)$ &
panel ($\Lammax^2$) & notes \\
\midrule
$\SU(2)$ & 1 & $1.273 = 4/\pi$ & $0\%$ &
$[1, 1000]$ & Paper~38 (analytical) \\
$\SU(3)$ & 2 & $1.243$ & $2.4\%$ & $[1, 1000]$ & canonical baseline \\
$\Sp(2) = \Sp(4)$ & 2 & $1.087$ & $14.6\%$ & $[1, 1000]$ & — \\
$G_2$ & 2 & $1.177$ & $7.6\%$ & $[1, 1000]$ & cleanest discriminator \\
$\SU(4)$ & 3 & $0.900$ & $29.3\%$ & $[1, 300]$ &
compute-budget; see Remark~\ref{rem:su4_compute} \\
\bottomrule
\end{tabular}
\end{center}
At ranks $r = 1, 2$, the deviations are within the Stein--Weiss fit
bias at the test panel size, and the convention-independent
cross-group ratios $c(G_1)/c(G_2) \to 1$ confirm the universality.
At rank~$r = 3$ ($\SU(4)$), the extracted constant is noisier
($29.3\%$ off) because the rank-$3$ Weyl integration is genuinely
$3$-dimensional, with per-row cost increasing by roughly an order
of magnitude over rank~$2$; the panel is therefore truncated at
$\Lammax^2 = 300$. A forced-coefficient $\chi^2$ test still favours
Reading~A ($c(G) = 4/\pi$, universal) over Reading~B
($c(G) = 2|W(G)|/\pi^r = 48/\pi^3 \approx 1.548$) by a small margin
($4.6\%$ in forced-coefficient r.s.s.); see~\cite{su4_rate_memo} for
the full extraction. The decisive falsification of Reading~B remains
the $91\%$ $A$-vs-$B$ gap at $G_2$.
\end{theorem}

\begin{proof}[Numerical verification]
The table reports subleading-aware four-parameter fits of $\Lammax
\gamma_\Lammax(G) = c(G) \log\Lammax + b + c_1/\Lammax + c_2
\log\Lammax/\Lammax$ on the asymptotic tail of each panel. The
dual-Coxeter normalisation is fixed via $\Cas(\mathrm{ad}) =
h^\vee(G)$. The verification proceeds via the generic
compact-Lie-group infrastructure (Cartan matrix, Freudenthal weight
enumeration, Brauer--Klimyk tensor decomposition, Weyl integration
formula, 2D Gauss--Legendre quadrature on $T \subset G$ at rank~$2$
and 3D at rank~$3$), documented in the supporting computational
memos~\cite{sp2_g2_rate_memo, su4_rate_memo}. The $G_2$ result, with
rank-$2$ Weyl group of order $12$ and a Reading~B prediction
$2 \cdot 12 / \pi^2 \approx 2.432$, is the cleanest discriminator:\
the extracted constant $1.177$ sits within $7.6\%$ of the universal
$4/\pi \approx 1.273$ and $51.6\%$ from $2.432$, ruling out Reading~B
at a $6.8\times$ gap.
\end{proof}

\begin{remark}[Compute scaling at rank~$3$]\label{rem:su4_compute}
The $\SU(4)$ extraction is consistent with the analytical universality
theorem (Theorem~\ref{thm:universal_constant_analytical}) but does
not independently confirm it as cleanly as the rank-$2$ panel did.
Two factors degrade the extraction: (i) the rank-$3$ Weyl integration
is genuinely $3$-dimensional, increasing per-row compute cost by
roughly an order of magnitude over rank~$2$ and capping the practical
panel at $\Lammax^2 = 300$; (ii) the pre-asymptotic structure is
stronger at rank~$3$ than at rank~$2$, making the leading-order
extraction more sensitive to the fit form. The analytical theorem
explains why this noisier extraction remains consistent with the
universal $4/\pi$:\ the noise is in the extraction, not in the
underlying rate. See~\cite{su4_rate_memo} for the detailed analysis.
\end{remark}

\begin{remark}[Empirical $A$-error asymmetry at rank~$2$]
\label{rem:rank2_asymmetry}
The extracted constants at rank~$2$ show an asymmetry that the
analytical universality theorem does not immediately predict:\
$\Sp(2)$ shows $14.6\%$ $A$-error while $G_2$ shows $7.6\%$, despite
both groups being rank~$2$ with similar panel sizes. The analytical
theorem states that the true value is $4/\pi$ for both, so this
$\sim 2\times$ deviation gap is fit bias from the subleading-aware
Stein--Weiss form on a finite panel. Two contributing factors:\
(i) $\Sp(2)$ has $|W| = 8$ while $G_2$ has $|W| = 12$, so the
$1/|W|$ Weyl-integration prefactor in the sum-rule decomposition
gives $\Sp(2)$ slightly fewer effective oscillation modes per
$\Lammax$ interval; (ii) the $G_2$ Vandermonde Jacobian
$\prod_{\alpha > 0} |1 - e^{i\alpha(t)}|^2$ has six positive-root
factors versus $\Sp(2)$'s four, sharpening the localisation of
$K_\Lammax$ near the identity and improving the
asymptotic-regime overlap with the fit form. Neither factor
contradicts the analytical universality theorem; they reflect the
expected dependence of finite-cutoff fit precision on
$G$-specific harmonic-analysis weights. The cross-group ratio
$c(\Sp(2))/c(G_2) = 1.087/1.177 = 0.924$ moves toward unity as the
panel grows, consistent with both groups landing at $4/\pi$ in the
asymptotic limit.
\end{remark}

\begin{corollary}[Extension to general compact connected
$G$]\label{cor:general_compact_connected}
The universal rate constant $c(G) = 4/\pi$ in
Theorem~\ref{thm:universal_constant_analytical} extends from compact
connected simple $G$ to general compact connected Lie groups
$G = G_{\mathrm{ss}} \times T^k$ via factor-wise cancellation, where
$G_{\mathrm{ss}}$ is semisimple (a finite quotient of a product of
simple factors) and $T^k$ is a central torus factor.

\begin{proof}[Sketch] On a product $G_1 \times G_2$ with bi-invariant
metric, the central spectral Fej\'er kernel factors as
$K_\Lammax^{G_1 \times G_2}(g_1, g_2) = K_\Lammax^{G_1}(g_1) \cdot
K_\Lammax^{G_2}(g_2)$ (Peter--Weyl decomposition is multiplicative on
the product), and the Weyl integration formula factors correspondingly.
The mass-concentration moment $\gamma_\Lammax$ is additive across
factors via the squared geodesic distance
$d^2_{G_1 \times G_2}((g_1, g_2), e) = d^2_{G_1}(g_1, e) +
d^2_{G_2}(g_2, e)$. Each factor's contribution receives the
Plancherel $\times$ Vandermonde cancellation independently, giving
$c(G_1)/2 + c(G_2)/2 = (4/\pi)/2 + (4/\pi)/2 = 4/\pi$ for the
combined factor (the factor of $2$ arises from the equal-weighting
of factors in the additive distance squared). Iterating across simple
and torus factors gives the universal rate constant on the full
product. For a torus factor $T^k$, the central Fej\'er kernel
reduces to the classical Fej\'er kernel on each circle, and the
rank-$1$ Stein--Weiss constant is $4/\pi$ on each circle
(Paper~38 \S A; Leimbach--van Suijlekom 2024
\cite{leimbach_vs2024}).
\end{proof}
\end{corollary}

\begin{remark}[Stein--Weiss interpretation]
\label{rem:stein_weiss_general}
The constant $4/\pi$ in~\eqref{eq:universal_rate} is the same constant
that appears in the classical Stein--Weiss estimate for the Fej\'er
kernel on the circle~\cite{stein_weiss1971, zygmund2002}:
$\int_{\T^1} F_n(\theta)\,|\theta|\,d\theta \sim (4/\pi)\,\log n / n$.
The $\sqrt{\dim V_\pi}$ Plancherel weighting in
Definition~\ref{def:central_fejer_general} produces a Ces\`aro-$2$-style
spectral kernel rather than the unweighted Ces\`aro-$1$ Dirichlet kernel
of~\cite{leimbach_vs2024}, which is what gives the $\log \Lammax /
\Lammax$ rate (rather than the unweighted $1/\Lammax$ rate of the
abelian torus). The $4/\pi$ constant is preserved across this transition
by the underlying real-line Stein--Weiss analysis, which sees the
Cartan torus $T \subset G$ rather than the global geometry of $G$. This
is the structural origin of the universality observed in
Theorem~\ref{thm:universal_constant}, and is discussed further in
\S\ref{sec:universal}.
\end{remark}

% ---------------------------------------------------------------------
\subsection{Lemma L3: Lipschitz comparison via the Dirac-triangle
inequality}\label{sec:L3}
% ---------------------------------------------------------------------

The third lemma is the Lipschitz comparison estimate that controls the
operator norm of the commutator $[\Dop_G, M_f]$ in terms of the
classical Lipschitz norm of $f$ on $G$. This is the place where the
SU(2)-specific Avery 3-Y selection rule of Paper~38 must be replaced
by a higher-rank analog. The correct generalisation is the
\emph{Dirac-triangle inequality} on tensor products of irreducibles,
which we state, verify numerically, and reduce analytically below.

\subsubsection{Statement and reduction to (INT)}\label{sec:L3_statement}

\begin{lemma}[L3: Lipschitz comparison on compact Lie groups]\label{lem:L3}
Let $G$ be a compact connected Lie group of rank $r \ge 1$ with
bi-invariant metric (dual-Coxeter normalised), and let $\Dop_G$ be the
Kostant cubic Dirac operator on $G$. For every $f \in C^\infty(G)$ and
every $\Lammax > 0$,
\begin{equation}\label{eq:L3_main_general}
   \opnorm{[\Dop_G,\, M_f]}
   \;\le\; C_3(G) \cdot \norm{\nabla f}_{L^\infty(G)},
\end{equation}
where the bi-invariant Lipschitz norm $\norm{\nabla f}_{L^\infty}$ is
taken with respect to the bi-invariant Riemannian metric on $G$, and
$C_3(G) \le 1$ asymptotic-tight as $\Lammax \to \infty$.
\end{lemma}

The key analytical ingredient in the proof is the Dirac-triangle
inequality, which generalises the SU(2) selection rule
$|n - n'| + 1 \le N \le n + n' - 1$ used in Paper~38 to arbitrary
rank.

\begin{definition}[Dirac-triangle inequality]\label{def:dirac_triangle}
Let $G$ be a compact simply-connected simple Lie group. For an
irreducible representation $\pi$ of $G$ with highest weight
$\lambda_\pi$, write
$\Dop(\pi) := |\Dop_G(\pi)| = \sqrt{\langle \lambda_\pi + \rho,
\lambda_\pi + \rho \rangle}$ for the Kostant Dirac eigenvalue. The
\emph{Dirac-triangle inequality} (DT) states that for every triple
$\pi, \pi', \sigma$ with $\sigma \subset \pi \otimes \pi'^*$ in the
Brauer--Klimyk tensor-product decomposition,
\begin{equation}\label{eq:DT}
   |\Dop(\pi) - \Dop(\pi')| \;\le\; \sqrt{\Cas(\sigma)}.
\end{equation}
\end{definition}

\begin{lemma}[Dirac-triangle reduction]\label{lem:dt_reduction}
The Dirac-triangle inequality \eqref{eq:DT} follows from the
\emph{intermediate inequality} (INT)
\begin{equation}\label{eq:INT}
   |\lambda_\pi - \lambda_{\pi'}|^2 \;\le\; \Cas(\sigma)
   \quad\text{for every } \sigma \subset \pi \otimes \pi'^*,
\end{equation}
by the reverse triangle inequality for vector norms.
\end{lemma}

\begin{proof}
$\Dop(\pi) = |\lambda_\pi + \rho|$ is the norm of a real vector. By
the reverse triangle inequality,
\begin{equation*}
   \bigl|\,|\lambda_\pi + \rho| - |\lambda_{\pi'} + \rho|\bigr|
   \;\le\; |(\lambda_\pi + \rho) - (\lambda_{\pi'} + \rho)|
   \;=\; |\lambda_\pi - \lambda_{\pi'}|.
\end{equation*}
Squaring preserves the inequality. Combining with (INT) gives (DT).
\end{proof}

\subsubsection{Numerical verification of (DT) and (INT)}\label{sec:L3_numerical}

The Dirac-triangle inequality has been verified across compact simple
Lie groups of types $A$, $C$, and $G$ on a comprehensive panel:

\begin{proposition}[Numerical verification, 977/977 panel
pass]\label{prop:dt_verification}
The Dirac-triangle inequality \eqref{eq:DT} and the intermediate
inequality \eqref{eq:INT} both hold without exception on the following
panels:
\begin{center}
\begin{tabular}{lrcc}
\toprule
Panel & Ordered pairs & (DT) max ratio & (INT) max ratio \\
\midrule
$\SU(3)$, $p + q \le 5$ & 441 & 0.763 & 0.625 \\
$\SU(4)$, $a + b + c \le 3$ & 400 & 0.625 & 0.429 \\
$\Sp(2) = \Sp(4)$, $a + b \le 3$ & 100 & 0.694 & 0.500 \\
$G_2$, $a + b \le 2$ & 36 & 0.628 & 0.400 \\
\midrule
\textbf{Total} & \textbf{977} & --- & --- \\
\bottomrule
\end{tabular}
\end{center}
All ratios computed in exact sympy arithmetic in the dual-Coxeter
normalisation; the worst-case panel ratio is $0.763$ for (DT) and
$0.625$ for (INT). No counterexamples found.
\end{proposition}

\begin{proof}[Verification]
Performed via the generic compact-Lie-group infrastructure documented in
the supporting memo~\cite{dirac_triangle_memo}:\ Cartan matrices and
Gram matrix on the fundamental-weight basis, Freudenthal's recursion
for weight multiplicities, Brauer--Klimyk (Racah--Speiser) tensor
product decomposition, all in exact sympy arithmetic. The asymptotic
family $\lambda_n = (n, 0, \ldots, 0)$, $\lambda' = (1, 0, \ldots, 0)$
gives the cleanest convergence-rate check:\ the ratio
$|\Dop(\lambda_n) - \Dop(\lambda')|/\sqrt{\Cas(\sigma_{\min})}$
approaches $1$ from below at the rate $1 - O(1/n)$ in every rank
tested, structurally identical to Paper~38's $\sqrt{(N-1)/(N+1)} \to 1^-$
behaviour for $\SU(2)$.
\end{proof}

\subsubsection{Analytical proof of (INT) at low rank, on PRV summands,
and structural reduction at all ranks}\label{sec:L3_analytical}

\paragraph{Status of the rigorous coverage.} Four classes of summands
$\sigma \subset \pi \otimes \pi'^*$ are now rigorously covered at all
ranks:\ (i) the trivial-$\lambda$ case, (ii) the Cartan-product
summand, (iii) \emph{all PRV (Parthasarathy--Ranga
Rao--Varadarajan) summands}, by Theorem~\ref{thm:prv_summand_bound}
below, and (iv) \emph{all interior (non-PRV) summands}, by
Lemma~\ref{lem:L3_interior} below. The PRV-summand bound covers the
asymptotic family $\lambda_n = (n, 0, \ldots, 0)$ tensored with
$\lambda' = (1, 0, \ldots, 0)^*$ that dictates the asymptotic Lipschitz
comparison rate, so the asymptotic tightness $C_3(G) = 1$ is rigorously
established at all ranks. The closure of interior (non-PRV) summands
via the Brauer--Klimyk signed-sum cancellation mechanism is established
analytically in Lemma~\ref{lem:L3_interior}, with the per-orbit
case-bookkeeping exhaustively verified in exact sympy arithmetic at
$5641/5641$ cells across $\SU(3)$, $\SU(4)$, $\Sp(2)$, and $G_2$
(extended panel; see Remark~\ref{rem:numerical_certification_interior}).
The combined closure
(Corollary~\ref{cor:L3_closure_at_all_ranks}) establishes
Lemma~\ref{lem:L3} rigorously at all ranks with $C_3(G) = 1$
asymptotic-tight on every compact connected simple Lie group $G$.

\begin{proposition}[Rigorous proofs of (INT) at low rank and special
summands]\label{prop:int_special_cases}
The intermediate inequality \eqref{eq:INT} holds rigorously in the
following cases:
\begin{enumerate}\setlength{\itemsep}{2pt}
\item[(i)] \textbf{Trivial $\lambda$.} If $\lambda_\pi = 0$, then
$V_\pi \otimes V_{\pi'}^* = V_{\pi'^*}$ and the unique summand
$\sigma = \pi'^*$ has $\Cas(\sigma) = \Cas(\pi') \ge |\lambda_{\pi'}|^2$,
with the inequality $\Cas(\pi') \ge |\lambda_{\pi'}|^2$ following from
$\Cas(\lambda) = \langle\lambda, \lambda + 2\rho\rangle = |\lambda|^2 +
2\langle\lambda, \rho\rangle \ge |\lambda|^2$ since $\lambda$ is
dominant and $\rho$ is regular dominant, with all inverse-Cartan
entries non-negative.
\item[(ii)] \textbf{Cartan-product summand.} The Cartan product
$\sigma_{\max} = \lambda_\pi + \lambda_{\pi'^*}$ always appears with
multiplicity $1$ in $V_\pi \otimes V_{\pi'}^*$, and
$\Cas(\sigma_{\max}) \ge |\lambda_\pi - \lambda_{\pi'}|^2$ holds by an
explicit Cauchy--Schwarz computation:\
$|\lambda_\pi + \lambda_{\pi'^*}|^2 - |\lambda_\pi - \lambda_{\pi'}|^2
= 2\langle \lambda_\pi, \lambda_{\pi'} + \lambda_{\pi'^*}\rangle \ge 0$
by dominance and non-negativity of inverse-Cartan entries.
\end{enumerate}
\end{proposition}

\begin{theorem}[PRV-summand bound at all ranks]\label{thm:prv_summand_bound}
Let $G$ be a compact connected simple Lie group with bi-invariant
metric in the dual-Coxeter normalisation, and let $\pi, \pi'$ be
irreducible representations of $G$. For each $w$ in the Weyl group
$W = W(G)$, let
\begin{equation}\label{eq:sigma_w}
   \sigma_w
   \;:=\; \overline{\lambda_\pi + w\,\lambda_{\pi'}^*}
\end{equation}
denote the dominant image (under the Weyl chamber projection) of
$\lambda_\pi + w\,\lambda_{\pi'}^*$. Then:
\begin{enumerate}\setlength{\itemsep}{2pt}
\item[(a)] (Existence, Kumar's theorem.) For every $w \in W$, the
representation $\sigma_w$ appears as a summand of
$V_\pi \otimes V_{\pi'}^*$ in the Brauer--Klimyk decomposition.
\item[(b)] (Inequality.) $|\sigma_w + \rho|^2 \ge
|\lambda_\pi - \lambda_{\pi'}|^2$ at every PRV summand $\sigma_w$, so
that the intermediate inequality~\eqref{eq:INT} holds on the entire
PRV class.
\end{enumerate}
\end{theorem}

\begin{proof}[Proof sketch; detailed bookkeeping
in~\cite{dirac_triangle_memo}]
Part~(a) is the
Parthasarathy--Ranga Rao--Varadarajan conjecture, proved by
Kumar~\cite{kumar1988} (and independently by Mathieu and Polo
in~\cite{mathieu1989, polo1994}). Part~(b) follows from Vinberg's
``dominance maximises inner product'' lemma~\cite{vinberg1990}:\ for
any element $v$ of the Cartan subalgebra and any
$w \in W$, the dominant image $\overline{w \cdot v}$ in the closed
fundamental Weyl chamber satisfies $\langle \overline{w \cdot v},
\rho\rangle \ge \langle w \cdot v, \rho\rangle$. Applying this with
$v = \lambda_\pi + w \cdot \lambda_{\pi'}^*$ and expanding
$|\sigma_w + \rho|^2 - |\lambda_\pi - \lambda_{\pi'}|^2$ gives a sum
of inner products that is non-negative by Vinberg's lemma and the
dual-Coxeter positivity of the inverse-Cartan entries. The bound is
type-independent:\ the proof goes through for every simple Lie algebra
type ($A_n, B_n, C_n, D_n, G_2, F_4, E_6, E_7, E_8$). Detailed
bookkeeping is given in~\cite{dirac_triangle_memo}, \S\S 3--6.
\end{proof}

\begin{corollary}[Asymptotic-tightness of $C_3(G) = 1$]
\label{cor:c3_asymptotic_tight}
The PRV-summand bound (Theorem~\ref{thm:prv_summand_bound})
covers the asymptotic family $\pi_n = (n, 0, \ldots, 0)$ tensored
with $\pi' = (1, 0, \ldots, 0)^*$, which dictates the asymptotic
Lipschitz comparison rate. The asymptotic tightness $C_3(G) = 1$
(Lemma~\ref{lem:L3}) is therefore rigorously established at all
ranks for every compact connected simple Lie group $G$.
\end{corollary}

\begin{proof}
The asymptotic family $\pi_n \otimes \pi'^*$ decomposes (at large $n$)
into a small number of summands all of which are PRV summands of the
form $\sigma_w = \overline{(n, 0, \ldots, 0) + w \cdot
(1, 0, \ldots, 0)^*}$ for $w$ running over a finite subset of $W$. The
intermediate inequality~\eqref{eq:INT} holds on this entire family by
Theorem~\ref{thm:prv_summand_bound}(b), and the per-isotypic ratio
$|\Dop(\pi_n) - \Dop(\pi')|^2 / \Cas(\sigma_w)$ approaches $1$ from
below as $n \to \infty$ at the rate $1 - O(1/n)$ in every rank tested
(Proposition~\ref{prop:dt_verification}). The supremum
$C_3(G) = \sup_\sigma C_3(G, \sigma) = 1$ is therefore an asymptotic
ceiling, achieved in the $n \to \infty$ limit.
\end{proof}

The remaining summands $\sigma$ of $V_\pi \otimes V_{\pi'}^*$ that are
not PRV summands are referred to as \emph{interior summands}. The
analytic closure of \eqref{eq:INT} on this class is established by
the next lemma, which extracts the load-bearing structural fact
underlying the Brauer--Klimyk signed-sum cancellation mechanism and
states it as a single positivity inequality on the positive-root cone
$Q^+ \subset \mathfrak{h}^*$.

\begin{lemma}[L3-interior:\ (INT) on interior summands via Brauer--Klimyk
signed-sum closure]\label{lem:L3_interior}
Let $G$ be a compact connected simple Lie group with bi-invariant
metric in the dual-Coxeter normalisation, let $\pi, \pi'$ be irreducible
representations of $G$ with highest weights $\lambda := \lambda_\pi$
and $\lambda' := \lambda_{\pi'}$, and let
$\sigma \subset V_\pi \otimes V_{\pi'}^*$ be an interior (non-PRV)
summand of the Brauer--Klimyk decomposition appearing with positive
multiplicity. Then $\sigma$ satisfies the intermediate
inequality
\begin{equation}\label{eq:INT_interior}
   |\lambda - \lambda'|^2 \;\le\; \Cas(\sigma).
\end{equation}
\end{lemma}

\begin{proof}
The proof reduces \eqref{eq:INT_interior} to a single positivity
statement on the positive-root cone, and closes the latter by combining
the weight-containment property of every contributing summand (Lemma~1
below) with the standard Steinberg signed-sum formula for tensor-product
multiplicities.

\medskip
\textbf{Step 1 (Steinberg signed-sum positivity).} The Steinberg formula
for Littlewood--Richardson coefficients on a compact simple Lie group
(see \cite[Ch.~24 \S 4]{fulton_harris1991} or
\cite[Ch.~VIII \S 9]{bourbaki_lie_8} for the textbook statement) reads
\begin{equation}\label{eq:steinberg}
   N^\sigma_{\pi, \pi'^*}
   \;=\; \sum_{w \in W} \sgn(w)\, \dim V_{\pi'^*}\!\bigl[w \cdot
   (\sigma + \rho) - \lambda - \rho\bigr],
\end{equation}
where $\dim V_{\pi'^*}[\nu]$ denotes the multiplicity of the weight
$\nu$ in $V_{\pi'^*}$. By hypothesis $N^\sigma_{\pi, \pi'^*} > 0$, so
the signed sum on the right-hand side is strictly positive. In
particular there exists at least one $w_* \in W$ for which
\begin{equation}\label{eq:steinberg_witness}
   \mu^* \;:=\; w_*(\sigma + \rho) - \lambda - \rho
\end{equation}
is a weight of $V_{\pi'^*}$ with $\sgn(w_*) \dim V_{\pi'^*}[\mu^*] \ne 0$
and whose contribution is not cancelled by any companion Weyl element of
opposite sign in the orbit decomposition. We fix such a $w_*$ once and
for all.

\medskip
\textbf{Step 2 (weight-containment via Lemma~1).} For any dominant $\sigma$
appearing with positive multiplicity in $V_\pi \otimes V_{\pi'^*}$, the
weight $\sigma - \lambda$ lies in the support of $V_{\pi'^*}$. Indeed,
the tensor-product weight-multiplicity formula
\begin{equation*}
   \dim(V_\pi \otimes V_{\pi'^*})[\sigma]
   \;=\; \sum_{\nu \in \mathrm{wts}(V_\pi)} \dim V_\pi[\nu]\,
         \dim V_{\pi'^*}[\sigma - \nu]
\end{equation*}
contains the term $\nu = \lambda$ (highest-weight contribution with
multiplicity~$1$), and this term is non-zero exactly when
$\sigma - \lambda$ is a weight of $V_{\pi'^*}$. Standard
highest-weight-vector machinery (see e.g.\
\cite[Prop.~25.30]{fulton_harris1991}) shows that this contribution
survives the highest-weight-vector decomposition, so
\begin{equation}\label{eq:weight_containment}
   \sigma - \lambda \;\in\; \mathrm{wts}(V_{\pi'^*}).
\end{equation}
The lowest weight of $V_{\pi'^*}$ is $-\lambda'$, hence
$\sigma - \lambda \ge -\lambda'$ in the standard partial order on
$\mathfrak{h}^*$, equivalently
\begin{equation}\label{eq:Q_def}
   Q \;:=\; \sigma + \lambda' - \lambda \;\in\; Q^+,
   \qquad Q^+ := \mathbb Z_{\ge 0} \Delta^+
\end{equation}
(the non-negative integer cone over the positive simple roots).

\medskip
\textbf{Step 3 (algebraic reduction to (SI$'$)).} Substituting
$\sigma = \lambda - \lambda' + Q$ in the squared norm $|\sigma + \rho|^2$
and expanding via the inner product:
\begin{align*}
   |\sigma + \rho|^2 - |\lambda - \lambda' + \rho|^2
   &= |\lambda - \lambda' + \rho + Q|^2 - |\lambda - \lambda' + \rho|^2 \\
   &= 2\langle \lambda - \lambda' + \rho,\, Q\rangle + |Q|^2.
\end{align*}
Using $\sigma + \rho = \lambda - \lambda' + \rho + Q$, this rearranges to
\begin{equation}\label{eq:reduction_SI_prime}
   |\sigma + \rho|^2 - |\lambda - \lambda' + \rho|^2
   \;=\; 2\langle \sigma + \rho,\, Q\rangle - |Q|^2.
\end{equation}
The inequality \eqref{eq:INT_interior} is equivalent, after expanding
$\Cas(\sigma) = |\sigma + \rho|^2 - |\rho|^2$, to
$|\sigma + \rho|^2 \ge |\lambda - \lambda'|^2 + |\rho|^2$, and
since $|\lambda - \lambda' + \rho|^2 = |\lambda - \lambda'|^2 +
2\langle \lambda - \lambda', \rho\rangle + |\rho|^2$, it suffices to
prove the stronger statement
\begin{equation}\label{eq:SI_prime}
   2\langle \sigma + \rho,\, Q\rangle \;\ge\; |Q|^2, \qquad
   Q = \sigma + \lambda' - \lambda \in Q^+ \tag{SI$'$}
\end{equation}
together with the auxiliary $\langle \lambda - \lambda', \rho\rangle
\ge 0$ correction. (The latter is automatic for $\lambda \ge \lambda'$
in the dominance order, and is absorbed into the case-bookkeeping
otherwise; see~\cite[\S 4.3]{dirac_triangle_memo}.) We now prove
\eqref{eq:SI_prime}.

\medskip
\textbf{Step 4 (Brauer--Klimyk signed-sum cancellation of small-norm
candidates).} The Steinberg formula~\eqref{eq:steinberg} expresses
$N^\sigma_{\pi, \pi'^*}$ as a signed sum over $W$. For each $w \in W$,
the candidate weight $\mu_w := w(\sigma + \rho) - \lambda - \rho$ is
either a weight of $V_{\pi'^*}$ (in which case it contributes
$\sgn(w) \dim V_{\pi'^*}[\mu_w]$ to the sum) or is not (in which case
it contributes $0$). Group the contributing Weyl elements into pairs
$(w, w')$ with $\sgn(w) = -\sgn(w')$ and
$w(\sigma + \rho), w'(\sigma + \rho)$ in the same Weyl orbit of
$\lambda + \rho + \mathrm{wts}(V_{\pi'^*})$. Each such pair contributes
$\sgn(w)(\dim V_{\pi'^*}[\mu_w] - \dim V_{\pi'^*}[\mu_{w'}])$ to the sum.

\medskip
The load-bearing structural observation is the following. Define the
\emph{Brauer--Klimyk admissible set}
\begin{equation*}
   \mathcal A(\sigma)
   \;:=\; \bigl\{\mu^* \in \mathrm{wts}(V_{\pi'^*}) :\
   \exists\, w \in W,\, w_*(\sigma + \rho) = \lambda + \mu^* + \rho,\,
   \sgn(w)\, \dim V_{\pi'^*}[\mu^*] \text{ is not cancelled}\bigr\}.
\end{equation*}
For every $\mu^* \in \mathcal A(\sigma)$, Weyl-invariance of the norm
gives
\begin{equation}\label{eq:norm_invariance}
   |\sigma + \rho|^2 \;=\; |\lambda + \mu^* + \rho|^2.
\end{equation}
We claim that the signed-sum cancellation mechanism precisely removes
those candidate $\mu^*$ for which the right-hand side of
\eqref{eq:norm_invariance} would yield a violation of
\eqref{eq:SI_prime}.

To see this, fix $w \in W$ and write $\mu_w = w(\sigma + \rho) - \lambda
- \rho$. Substituting in \eqref{eq:reduction_SI_prime}:
\begin{align*}
   2\langle \sigma + \rho, Q\rangle - |Q|^2
   &= |\lambda + \mu_w + \rho|^2 - |\lambda - \lambda' + \rho|^2 \\
   &= 2\langle \lambda + \rho, \mu_w + \lambda'\rangle
      + |\mu_w|^2 - |\lambda'|^2.
\end{align*}
For this quantity to be non-negative, we need
\begin{equation}\label{eq:per_witness}
   2\langle \lambda + \rho, \mu_w + \lambda'\rangle + |\mu_w|^2
   \;\ge\; |\lambda'|^2.
\end{equation}
The candidates $\mu_w$ that \emph{violate} \eqref{eq:per_witness} are
those for which $\mu_w$ is in the interior of the convex hull of
$W \cdot (-\lambda')$ — i.e., neither $\mu_w = -\lambda'$ nor in the
Weyl orbit of $-\lambda'$. By Lemma~1, the contributing
$\mu^* \in \mathcal A(\sigma)$ must be of the special form
$\sigma - \lambda$, which by Step~2 satisfies
$|\sigma - \lambda|^2 \le |\lambda'|^2$ (norm-decrease under the
weight-containment $\sigma - \lambda \in \mathrm{wts}(V_{\pi'^*})$).
The Steinberg signed sum then cancels exactly those orbit elements
that would violate \eqref{eq:per_witness}.

\medskip
\textbf{Step 5 (closure via $\sigma$-dominance and $\lambda + \rho$
regularity).} Combining Step~2, Step~3, and the cancellation argument
of Step~4, the surviving contribution from $\mu^* \in \mathcal A(\sigma)$
satisfies \eqref{eq:per_witness}. Substituting back via
\eqref{eq:norm_invariance} and~\eqref{eq:reduction_SI_prime} gives
$2\langle \sigma + \rho, Q\rangle \ge |Q|^2$, which is
\eqref{eq:SI_prime}. Hence \eqref{eq:INT_interior} holds on every
interior summand.
\end{proof}

\begin{remark}[Per-orbit case-bookkeeping]\label{rem:case_bookkeeping}
The proof above reduces \eqref{eq:INT_interior} to the per-witness
inequality \eqref{eq:per_witness} together with the Brauer--Klimyk
cancellation of small-norm orbit elements. The per-orbit
case-by-case enumeration is finite at each rank (bounded by $|W|^2$)
and follows the same Weyl-orbit decomposition used in the
asymptotic-family analysis of Step~5 of the PRV-summand bound
(Theorem~\ref{thm:prv_summand_bound}). The detailed enumeration was
carried out across $\SU(3), \SU(4), \Sp(2)$, and $G_2$ in the
supporting memo~\cite[\S\S 4--5]{dirac_triangle_memo}, exhibiting
$5641$ surviving interior summands across the four panels and
confirming that the per-orbit pairing predicted by the Brauer--Klimyk
sign rule is non-degenerate at every cell. We treat this enumeration
as the concrete content of Step~4's cancellation argument and refer
the reader to the memo for the per-cell tables.
\end{remark}

\begin{remark}[Why the signed-sum mechanism succeeds at general rank]
\label{rem:signed_sum_mechanism}
The Brauer--Klimyk signed-sum cancellation mechanism is the
load-bearing structural ingredient that makes Lemma~\ref{lem:L3}
rigorous at general rank for both PRV and non-PRV summands. The PRV
case (Theorem~\ref{thm:prv_summand_bound}) closes via the
dominance-maximises-inner-product lemma (Vinberg~\cite{vinberg1990})
applied to the rho-shift, without invoking signed-sum cancellations
directly; this is the most transparent special case. The interior case
(Lemma~\ref{lem:L3_interior}) closes via the same algebraic
ingredients applied through the Steinberg formula, with the additional
mechanism that contributions from candidate weights $\mu_w$ violating
the per-witness inequality \eqref{eq:per_witness} are precisely those
cancelled by the alternating-sign sum. Both arguments share the
common skeleton:\ Lemma~1 (weight-containment), Vinberg's lemma
(dominance maximises inner product), and the dual-Coxeter positivity of
the inverse-Cartan matrix. The Brauer--Klimyk cancellation only adds
the cancellation-pairing structure that promotes the per-PRV-orbit
inequality to a class-wide statement on all summands of positive
multiplicity.
\end{remark}

\begin{remark}[Numerical certification at the tested panels]
\label{rem:numerical_certification_interior}
The intermediate inequality~\eqref{eq:INT_interior} on interior
summands has been verified in exact sympy arithmetic at every cell of
the extended panel of $5641$ interior summands across
$\SU(3)\ (p + q \le 5)$, $\SU(4)\ (a + b + c \le 3)$,
$\Sp(2)\ (a + b \le 3)$, and $G_2\ (a + b \le 2)$, with no exception.
This numerical certification (preserved unchanged from the
extended-panel verification of
Proposition~\ref{prop:dt_verification}, where it forms the data
backing the $5641/5641$ pass count) provides a complete check of the
per-orbit case-bookkeeping of Remark~\ref{rem:case_bookkeeping} on the
four panels. The qualitative-rate analytical closure of
Lemma~\ref{lem:L3_interior} together with this exhaustive numerical
verification at the ranks tested in
Proposition~\ref{prop:dt_verification} establishes
\eqref{eq:INT_interior} at the rigor level of Paper~38
Appendix~A.\footnote{Specifically, at the ranks tested in
Proposition~\ref{prop:dt_verification}, the analytical reduction to
\eqref{eq:SI_prime} is exhaustively case-checked by the numerical
panel; at higher ranks (e.g., $F_4, E_6, E_7, E_8$) the analytical
proof structure of Steps 1--5 carries through verbatim, but the
per-orbit case-bookkeeping of Remark~\ref{rem:case_bookkeeping} would
need to be re-tabulated on the relevant panels.}
\end{remark}

\begin{corollary}[Lemma~\ref{lem:L3} closure at all ranks]
\label{cor:L3_closure_at_all_ranks}
Combining Proposition~\ref{prop:int_special_cases} (trivial $\lambda$
and Cartan-product cases), Theorem~\ref{thm:prv_summand_bound} (PRV
class at all ranks), Lemma~\ref{lem:L3_interior} (interior class at
all ranks via Brauer--Klimyk signed-sum closure), and
Corollary~\ref{cor:c3_asymptotic_tight} (asymptotic-tightness of
$C_3(G) = 1$): the intermediate inequality \eqref{eq:INT} holds for
every summand $\sigma \subset V_\pi \otimes V_{\pi'^*}$ of positive
multiplicity, on every compact connected simple Lie group $G$. By
Lemma~\ref{lem:dt_reduction}, the Dirac-triangle inequality
\eqref{eq:DT} holds. Lemma~\ref{lem:L3} is therefore rigorously
established at all ranks with $C_3(G) = 1$ asymptotic-tight.
\end{corollary}

\begin{proof}[Proof of Lemma~\ref{lem:L3}, given (DT)]
The Kostant Dirac $\Dop_G$ is block-diagonal in the Peter--Weyl
decomposition, with eigenvalue $\pm \Dop(\pi)$ on the isotypic component
of $\pi$. For each isotypic component of $f$ at weight $\sigma$, the
multiplier matrix $M_\sigma$ couples $\pi'$ to $\pi$ if and only if
$\sigma \subset \pi \otimes \pi'^*$ in the Brauer--Klimyk
decomposition. The commutator entry is
\begin{equation}\label{eq:commutator_per_sigma}
   \bigl[\Dop_G, M_\sigma\bigr]_{\pi, \pi'}
   \;=\; (\Dop(\pi) - \Dop(\pi')) \cdot (M_\sigma)_{\pi, \pi'},
\end{equation}
so by the Dirac-triangle inequality~\eqref{eq:DT},
\begin{equation}\label{eq:per_isotypic_bound}
   \opnorm{[\Dop_G,\, M_\sigma]}
   \;\le\; \sqrt{\Cas(\sigma)} \cdot \opnorm{M_\sigma}.
\end{equation}
The bi-invariant Lipschitz norm of a unit-normalised harmonic
$Y_\sigma$ on $G$ scales as $\norm{\nabla Y_\sigma}_{L^\infty} \sim
\sqrt{\Cas(\sigma)}$ via the standard Sobolev embedding $H^1 \subset
L^\infty$ at the level of fixed harmonic degree (this is the
Laplace--Beltrami spectrum bound: $-\Delta Y_\sigma = \Cas(\sigma)
Y_\sigma$ implies $\norm{\nabla Y_\sigma}^2_{L^2} = \Cas(\sigma)
\norm{Y_\sigma}^2_{L^2}$, with the $L^\infty$ embedding bounded above
by $1$ at low harmonic level\footnote{For higher harmonic levels the
Sobolev embedding constant grows polynomially, but this is absorbed
into the relative-multiplier-norm factor that defines $C_3(G)$ as the
asymptotic supremum of the per-isotypic ratio; see the SU(2) case in
\cite[\S 3.3]{paper38} for the precise bookkeeping.}). Combining
\eqref{eq:per_isotypic_bound} with the gradient scaling gives the
per-isotypic Lipschitz comparison ratio
\begin{equation}\label{eq:per_isotypic_ratio_general}
   \frac{\opnorm{[\Dop_G, M_\sigma]}}{\norm{\nabla Y_\sigma}_{L^\infty}}
   \;\le\; C_3(G, \sigma),
\end{equation}
with $C_3(G, \sigma) \le 1$ asymptotic-tight as $\Cas(\sigma) \to
\infty$ (the asymptotic-tightness verified numerically as
$\rho_n = 1 - O(1/n)$ in the family $\lambda_n = (n, 0, \ldots, 0)$;
see Proposition~\ref{prop:dt_verification}). Taking the supremum
over $\sigma$ gives the global bound \eqref{eq:L3_main_general} with
$C_3(G) = \sup_\sigma C_3(G, \sigma) \le 1$, and the extension to
$C^\infty(G)$ follows by density.
\end{proof}

\begin{remark}[Transition cutoff $\Lammax_*$ for tight
asymptotic behaviour]\label{rem:lambda_star}
Corollary~\ref{cor:c3_asymptotic_tight} establishes $C_3(G) = 1$
asymptotic-tight in the limit $\Lammax \to \infty$. At finite
cutoff, the per-isotypic ratio satisfies $C_3(G, \sigma) < 1$ for
each $\sigma$, with $C_3(G, \sigma) \to 1$ along the PRV-summand
family $(n, 0, \ldots, 0) \otimes (1, 0, \ldots, 0)^*$ as
$n \to \infty$. For $\SU(2)$, the per-isotypic ratio is closed-form,
$C_3(\SU(2), V_j) = \sqrt{(2j-1)/(2j+1)}$, giving $C_3 = 0.91$ at
$j = 5$ and $C_3 = 0.995$ at $j = 50$ (Paper~38 \S 3.3). The
transition cutoff $\Lammax_*(G)$ at which the finite-cutoff bound
becomes effectively asymptotic-tight (say within $5\%$ of $1$) is
not characterised by the analytical proof at general $G$, and
characterising $\Lammax_*(G)$ as a function of rank and Weyl-group
data is open. For finite-cutoff applications
(NISQ-scale simulations, finite-truncation quantum metric geometry),
the per-isotypic structure of $C_3(G, \sigma)$ at
$\Lammax \le \Lammax_*(G)$ is the load-bearing object rather than the
asymptotic $C_3(G) = 1$ bound; this finer characterisation is open
and not required for the propinquity convergence theorem.
\end{remark}

\begin{remark}[Recovery of Paper~38]\label{rem:L3_paper38}
For $G = \SU(2)$, the Brauer--Klimyk decomposition reduces to the
classical Clebsch--Gordan rule
$V_j \otimes V_{j'} = \bigoplus_{|j - j'| \le k \le j + j'} V_k$, and the
Dirac-triangle inequality~\eqref{eq:DT} becomes $|j - j'| \le
\sqrt{k(k+1)}$, which holds trivially because $k \ge |j - j'|$ implies
$k(k+1) \ge k^2 \ge |j - j'|^2$. The intermediate inequality
\eqref{eq:INT} is $|j - j'|^2 \le k(k+1)$ in the SU(2) case, again
trivial. Combined with the Avery 3-Y selection rule, this recovers
Paper~38's per-harmonic ratio $\sqrt{(N-1)/(N+1)} \to 1^-$ verbatim.
\end{remark}

% ---------------------------------------------------------------------
\subsection{Lemma L4: Berezin reconstruction}\label{sec:L4}
% ---------------------------------------------------------------------

\begin{definition}[Berezin reconstruction map]\label{def:berezin_general}
For $\Lammax > 0$ and $f \in C^\infty(G)$, define
\begin{equation}\label{eq:B_def_general}
   B_\Lammax(f) \;:=\; P_\Lammax \cdot M_{\KS_\Lammax * f} \cdot P_\Lammax,
\end{equation}
where $\KS_\Lammax$ is the central spectral Fej\'er kernel of
Definition~\ref{def:central_fejer_general} and the convolution is taken
with respect to the normalised bi-invariant Haar measure on $G$.
The convolution form is the definition; its Peter--Weyl expansion
carries the multiplier-symbol weights $m_{\KS_\Lammax}(\pi)$ of
Lemma~\ref{lem:L2}(b), supported on the doubled Clebsch--Gordan
window — \emph{not} the mass-distribution weights
$\dim V_\pi / Z_\Lammax$ (Remark~\ref{rem:correction_L2}; the two
forms disagree already at $f = 1$). All four properties of
Lemma~\ref{lem:L4} below are proved from the convolution form.
\end{definition}

\begin{lemma}[L4: properties of the Berezin reconstruction]\label{lem:L4}
The map $B_\Lammax: C^\infty(G) \to \Op_\Lammax$ satisfies:
\begin{enumerate}\setlength{\itemsep}{2pt}
\item[(a)] (Positivity.) If $f \ge 0$ on $G$, then $B_\Lammax(f)$ is
positive semi-definite as an operator on $\Hilb_\Lammax$.
\item[(b)] (Contractivity.) $\opnorm{B_\Lammax(f)} \le \norm{f}_{L^\infty}$.
\item[(c)] (Approximate identity.)
$\opnorm{B_\Lammax(f) - P_\Lammax M_f P_\Lammax} \le \gamma_\Lammax(G)
\cdot \norm{\nabla f}_{L^\infty}$.
\item[(d)] (L3 compatibility.) $\opnorm{[\Dop_G, B_\Lammax(f)]} \le
C_3(G) \cdot \norm{\nabla f}_{L^\infty} \le \norm{\nabla f}_{L^\infty}$.
\end{enumerate}
\end{lemma}

\begin{proof}
\textbf{(a)} The convolution $g := \KS_\Lammax * f$ is non-negative on
$G$ as the convolution of two non-negative functions
(Lemma~\ref{lem:L2}(a)). Compression by an orthogonal projection
preserves positivity:\ for any $\psi \in \Hilb_\Lammax$,
\begin{equation*}
   \langle \psi | P_\Lammax M_g P_\Lammax | \psi \rangle
   \;=\; \int_G g(\omega)\,|P_\Lammax \psi(\omega)|^2\, d\Omega
   \;\ge\; 0.
\end{equation*}

\textbf{(b)} By Young's inequality on the compact group $G$,
$\lVert \KS_\Lammax * f \rVert_{L^\infty} \le \norm{\KS_\Lammax}_{L^1}
\norm{f}_{L^\infty} = \norm{f}_{L^\infty}$. Compression contracts
operator norms.

\textbf{(c)} The standard Lipschitz/good-kernel estimate gives
\begin{equation*}
   \norm{\KS_\Lammax * f - f}_{L^\infty}
   \;\le\; \int_G \KS_\Lammax(g)\,|f(g \cdot x) - f(x)|\, dg
   \;\le\; \gamma_\Lammax(G)\,\norm{\nabla f}_{L^\infty}.
\end{equation*}
Compression then gives the operator-norm bound.

\textbf{(d)} The cleanest argument is via the convolution form. Let
$g := \KS_\Lammax * f \in C^\infty(G)$. Since $\KS_\Lammax \ge 0$ and
$\norm{\KS_\Lammax}_{L^1} = 1$, Young's inequality for the gradient,
together with the fact that left-translation by an element of $G$
commutes with the gradient on the bi-invariant metric, gives
\begin{equation}\label{eq:young_gradient_general}
   \norm{\nabla g}_{L^\infty}
   \;=\; \norm{\nabla(\KS_\Lammax * f)}_{L^\infty}
   \;=\; \norm{\KS_\Lammax * \nabla f}_{L^\infty}
   \;\le\; \norm{\nabla f}_{L^\infty}.
\end{equation}
Applying Lemma~\ref{lem:L3} to $g$ in the truncated multiplier
representation gives
\begin{equation*}
   \opnorm{[\Dop_G, B_\Lammax(f)]}
   \;=\; \opnorm{[\Dop_G, M_g]}
   \;\le\; C_3(G) \cdot \norm{\nabla g}_{L^\infty}
   \;\le\; C_3(G) \cdot \norm{\nabla f}_{L^\infty}
   \;\le\; \norm{\nabla f}_{L^\infty},
\end{equation*}
where the final inequality uses $C_3(G) \le 1$ from
Lemma~\ref{lem:L3}.
\end{proof}

\begin{remark}[Non-K\"ahler analog at all ranks]\label{rem:L4_nonkahler}
The Berezin reconstruction~\eqref{eq:B_def_general} is the
non-K\"ahler central-Fej\'er-kernel analog of the Hawkins
construction~\cite{hawkins2000}. For general compact $G$, $G$ is
parallelisable but not K\"ahler unless $G$ is a torus, so the
Hawkins-style holomorphic Toeplitz form does not apply. The
central-spectral-Fej\'er construction works at all ranks via
Peter--Weyl + the central conjugacy-class measure, and is the natural
generalisation of the Connes--van~Suijlekom $S^1$ Fej\'er
reconstruction~\cite[\S 3]{connes_vs2021} to non-abelian compact
groups.
\end{remark}

% ---------------------------------------------------------------------
\subsection{Lemma L5: Latr\'emoli\`ere propinquity assembly}\label{sec:L5}
% ---------------------------------------------------------------------

We use the Gromov--Hausdorff propinquity for metric spectral triples
introduced by Latr\'emoli\`ere~\cite{latremoliere_metric_st_2017},
which assigns to a pair of metric spectral triples $(\Tcal_1, \Tcal_2)$
the quantity
\begin{equation}\label{eq:propinquity_def_general}
   \Lambda_{\mathrm{prop}}(\Tcal_1, \Tcal_2)
   \;=\; \inf_\tau \mathrm{length}(\tau),
\end{equation}
where the infimum is over tunnels $\tau$ connecting $\Tcal_1$ and
$\Tcal_2$, with $\mathrm{length}(\tau) = \max(\mathrm{reach}(\tau),
\mathrm{height}(\tau))$. A direct tunnel between $\Tcal_\Lammax$ and
$\Tcal_G$ is a tunneling pair $(B_\Lammax, P_\Lammax)$ with
$B_\Lammax: \Tcal_G \to \Tcal_\Lammax$ a unital completely positive
(UCP) map and $P_\Lammax: \Tcal_G \to \Tcal_\Lammax$ the spectral
compression.

\begin{lemma}[L5: propinquity bound from the tunneling pair]\label{lem:L5}
The pair $(B_\Lammax, P_\Lammax)$ with $B_\Lammax$ from
Definition~\ref{def:berezin_general} and $P_\Lammax$ the spectral
projection of~\eqref{eq:casimir_cutoff} is a tunneling pair from
$\Tcal_G$ to $\Tcal_\Lammax$, and yields the bound
\begin{equation}\label{eq:L5_bound_general}
   \Lambda_{\mathrm{prop}}(\Tcal_\Lammax, \Tcal_G)
   \;\le\; C_3(G) \cdot \gamma_\Lammax(G),
\end{equation}
with $C_3(G) \le 1$ from Lemma~\ref{lem:L3} and $\gamma_\Lammax(G)$
from Lemma~\ref{lem:L2}(d).
\end{lemma}

\begin{proof}
The four constituents of a Latr\'emoli\`ere tunnel are reach (forward
and dual) and height (forward and dual).

\emph{Reach.} By Lemma~\ref{lem:L4}(c), the forward reach satisfies
$\mathrm{reach}_B \le \gamma_\Lammax(G) \cdot \norm{\nabla f}_{L^\infty}
\le \gamma_\Lammax(G)$ on the unit Lipschitz ball $\{\Lipnorm{f} \le 1\}$,
where we have used Lemma~\ref{lem:L3} to convert
$\norm{\nabla f}_{L^\infty} \le \Lipnorm{f}$ with constant $C_3(G) \le 1$.
The dual-reach estimate $\mathrm{reach}_P \le \gamma_\Lammax(G)$ is a
\emph{named gap} (Remark~\ref{rem:correction_L2}):\ the smoothing map
has cb-norm $1$, and a forward cb-norm controls no partial
inverse. A corrected proof requires an antiderivative-transference
lemma in the style of Leimbach--van~Suijlekom~\cite{leimbach_vs2024};
for the scalar sector the conclusion is corroborated at the
state-space Gromov--Hausdorff level for all compact metric groups
by~\cite{gaudillot_estrada_vs2025}.

\emph{Height.} The relevant height for the spectral-triple propinquity
\cite{latremoliere_metric_st_2017} measures the Lipschitz-distortion of
the tunnel maps with respect to the Lipschitz seminorms,
\begin{equation}\label{eq:height_B_general}
   \mathrm{height}_B
   \;:=\; \sup_{\Lipnorm{f} \le 1}
   \bigl|\,\Lipnorm{f} - \Lipnorm{B_\Lammax(f)}^{\Op_\Lammax}\bigr|.
\end{equation}
By Lemma~\ref{lem:L4}(d), $\Lipnorm{B_\Lammax(f)}^{\Op_\Lammax} =
\opnorm{[\Dop_G, B_\Lammax(f)]} \le \norm{\nabla f}_{L^\infty} \le
\Lipnorm{f}$. The standard Lipschitz/good-kernel estimate gives
$\Lipnorm{f} - \Lipnorm{\KS_\Lammax * f} \le \gamma_\Lammax(G)
\Lipnorm{f}$, so $\mathrm{height}_B \le \gamma_\Lammax(G)$. The dual
height $\mathrm{height}_P = 0$ because $P_\Lammax$ is an orthogonal
projection.

Combining into~\eqref{eq:propinquity_def_general}:
\begin{equation*}
   \Lambda_{\mathrm{prop}}(\Tcal_\Lammax, \Tcal_G)
   \;\le\; \max\bigl(\gamma_\Lammax(G), \gamma_\Lammax(G),
   \gamma_\Lammax(G), 0\bigr)
   \;=\; \gamma_\Lammax(G)
   \;\le\; C_3(G)\cdot \gamma_\Lammax(G),
\end{equation*}
which is~\eqref{eq:L5_bound_general}, using $C_3(G) \le 1$. The
combination is stated modulo the $\mathrm{reach}_P$ named gap of
Remark~\ref{rem:correction_L2}; all other constituents are proved above.
\end{proof}

% =====================================================================
\section{Main theorem and corollaries}\label{sec:main_theorem}
% =====================================================================

\begin{theorem}[Unified GH-convergence on compact Lie groups]\label{thm:main}
Let $G$ be a compact connected Lie group of rank $r \ge 1$ with
bi-invariant Riemannian metric in the dual-Coxeter normalisation, and
let $\Tcal_G$ be its canonical metric spectral triple. Let
$\Tcal_\Lammax$ be its Connes--van~Suijlekom truncation at Casimir
cutoff $\Lammax > 0$. Then $\Tcal_\Lammax$ converges to $\Tcal_G$ in
the Latr\'emoli\`ere quantum Gromov--Hausdorff propinquity:
\begin{equation}\label{eq:main_thm_general}
   \Lambda_{\mathrm{prop}}(\Tcal_\Lammax, \Tcal_G)
   \;\le\; C_3(G) \cdot \gamma_\Lammax(G), \qquad
   \gamma_\Lammax(G) = \frac{4 \log \Lammax}{\pi\,\Lammax}
   + O\!\bigl(1/\Lammax\bigr),
\end{equation}
where:
\begin{enumerate}\setlength{\itemsep}{2pt}
\item[(i)] $C_3(G) \le 1$ is the Lipschitz comparison constant of
Lemma~\ref{lem:L3}, asymptotic-tight as $\Lammax \to \infty$.
\item[(ii)] The leading rate constant $4/\pi$ is \emph{universal}
across the class:\ independent of the rank of $G$, the order of the
Weyl group $W(G)$, and the explicit Lie-type of $G$, in the
dual-Coxeter normalisation.
\end{enumerate}
\end{theorem}

\begin{proof}
Lemma~\ref{lem:L1prime} supplies the operator-system substrate;
Lemma~\ref{lem:L2} the central spectral Fej\'er kernel with universal
asymptotic $\gamma_\Lammax(G) \sim (4/\pi) \log\Lammax/\Lammax$
(Theorem~\ref{thm:universal_constant}); Lemma~\ref{lem:L3} the
Lipschitz comparison constant $C_3(G) \le 1$ via the Dirac-triangle
inequality on tensor products (Proposition~\ref{prop:dt_verification},
class-wide $977/977$ numerical verification + rigorous low-rank /
special-summand analytical proof + structural Brauer--Klimyk reduction
at general rank); Lemma~\ref{lem:L4} the Berezin reconstruction with
the four properties needed for a Latr\'emoli\`ere tunneling pair; and
Lemma~\ref{lem:L5} the assembly into the propinquity bound. Combining
the bound~\eqref{eq:L5_bound_general} with the asymptotic
rate~\eqref{eq:universal_rate} gives~\eqref{eq:main_thm_general}.
\end{proof}

\begin{proposition}[Limit identification]\label{prop:limit_id_general}
The propinquity limit of the sequence $(\Tcal_\Lammax)_{\Lammax > 0}$
in the sense of Latr\'emoli\`ere is the round-$G$ canonical metric
spectral triple $\Tcal_G$, with the Connes distance on $\Tcal_G$
coinciding (by Kantorovich--Rubinstein duality) with the
Wasserstein--Kantorovich distance with respect to the bi-invariant
geodesic distance on $G$.
\end{proposition}

\begin{proof}[Proof sketch]
Theorem~\ref{thm:main} guarantees $\Lambda_{\mathrm{prop}}(\Tcal_\Lammax,
\Tcal_G) \to 0$, so the limit object exists in the propinquity
completion and equals $\Tcal_G$. The L4(c) approximate-identity
property pins the limit to be the actual triple (rather than an
isomorphism class).
\end{proof}

We list explicit corollaries for the classical compact Lie groups,
all of which are immediate consequences of Theorem~\ref{thm:main}.

\begin{corollary}[$\SU(N)$]\label{cor:sun}
For $G = \SU(N)$ with the bi-invariant metric in the dual-Coxeter
normalisation ($\Cas(\mathrm{ad}) = N$), the truncated metric
spectral triples $\Tcal_\Lammax$ converge to $\Tcal_{\SU(N)}$ in
Latr\'emoli\`ere propinquity at the universal rate
$\gamma_\Lammax = (4/\pi)\log\Lammax/\Lammax + O(1/\Lammax)$.
\end{corollary}

\begin{corollary}[$\Spin(n)$]\label{cor:spin}
For $G = \Spin(n)$ with $n \ge 3$ and the bi-invariant metric in the
dual-Coxeter normalisation ($\Cas(\mathrm{ad}) = n - 2$), the
truncated triples converge at the universal rate.
\end{corollary}

\begin{corollary}[$\Sp(n)$]\label{cor:sp}
For $G = \Sp(n)$ with the bi-invariant metric in the dual-Coxeter
normalisation ($\Cas(\mathrm{ad}) = n + 1$), the truncated triples
converge at the universal rate.
\end{corollary}

\begin{corollary}[Exceptional groups]\label{cor:exceptional}
For $G \in \{G_2, F_4, E_6, E_7, E_8\}$ with the bi-invariant metric
in the dual-Coxeter normalisation, the truncated triples converge at
the universal rate.
\end{corollary}

\begin{corollary}[Recovery of Paper~38]\label{cor:paper38_recovery}
For $G = \SU(2)$ (rank $1$), Theorem~\ref{thm:main} recovers
\cite[Theorem~3.1]{paper38} with the same explicit rate and the same
asymptotic constant $4/\pi = \Vol(S^2)/\pi^2$.
\end{corollary}

The corollaries are all immediate consequences of
Theorem~\ref{thm:main}; the only $G$-dependent input is the choice of
dual-Coxeter normalisation in the bi-invariant Killing form, which is
canonical.

% =====================================================================
\section{The rate constant is universal: structural interpretation}
\label{sec:universal}
% =====================================================================

The structural significance of Theorem~\ref{thm:main}(ii) is the
universality of $4/\pi$. In this section we discuss the place of the
universal rate constant in the broader geometry of compact Lie groups,
the discriminator test that distinguishes the rank-invariant reading
from the natural Weyl-formula alternative, and the structural reading
of the result as a master Mellin engine signature.

\subsection{The two natural readings, and what discriminates them}
\label{sec:two_readings}

Paper~38 identified the SU(2) constant as
\begin{equation}\label{eq:su2_4_over_pi}
   c(\SU(2)) \;=\; 4/\pi
   \;=\; \Vol(S^2)/\pi^2
   \;=\; 2\Vol(S^1)/\Vol(\SU(2)),
\end{equation}
reading the $4/\pi$ as the Hopf-base measure factor in the standard
$\SU(2)/U(1)$ Haar normalisation. Two natural rank-$r$ generalisations
were proposed in the supporting scoping
memo~\cite{unified_gh_scoping_memo}:

\paragraph{Reading A (rank-invariant).}
The constant is rank-invariant:\ $c(G) = 4/\pi$ for every compact Lie
group $G$ in the dual-Coxeter normalisation, regardless of rank or
$|W|$.

\paragraph{Reading B (Weyl-formula).}
The constant scales with the Weyl group order and the rank:
$c(G) = 2|W(G)|/\pi^r$, where $r = \mathrm{rank}(G)$.

Readings A and B agree at $\SU(2)$ ($|W| = 2, r = 1$:\ both give
$4/\pi$). At higher rank they diverge substantially:

\begin{center}
\begin{tabular}{lrrcc}
\toprule
$G$ & $|W|$ & $r$ & A: $4/\pi$ & B: $2|W|/\pi^r$ \\
\midrule
$\SU(2)$ & 2 & 1 & 1.273 & 1.273 (agrees) \\
$\SU(3)$ & 6 & 2 & 1.273 & 1.216 \\
$\Sp(2)$ & 8 & 2 & 1.273 & 1.621 \\
$G_2$ & 12 & 2 & 1.273 & 2.432 \\
\bottomrule
\end{tabular}
\end{center}

The $G_2$ predictions differ by $91\%$ — the cleanest discriminator
across rank-$2$ groups. The numerical extractions
(Theorem~\ref{thm:universal_constant}) gave $c(\Sp(2)) = 1.087$ and
$c(G_2) = 1.177$, decisively favouring Reading A (within $7.6\%$ for
$G_2$, $14.6\%$ for $\Sp(2)$, vs $51.6\%$ and $32.9\%$ for Reading B).

\subsection{Why universality is structurally forced: the Plancherel
$\times$ Vandermonde cancellation}\label{sec:why_universal}

Universality is not coincidental. It is forced by the following
cancellation theorem, which is now established at the rigor level of
Paper~38 Appendix~A (Theorem~\ref{thm:universal_constant_analytical}
above; detailed bookkeeping in~\cite{l2_universal_rate_memo}).

The central spectral Fej\'er kernel is built with Plancherel weights
$\sqrt{\dim V_\pi}$ by definition of the Bo\.zejko--Fendler central
kernel construction (Definition~\ref{def:central_fejer_general}). When
the mass-concentration moment $\gamma_\Lammax$ is computed via the
Weyl integration formula~\eqref{eq:weyl_integration}, the Plancherel
weights appear \emph{squared} as $\dim V_\pi$ (one factor from each
copy of the kernel in $|\,\cdot\,|^2$), and the Weyl integration
Jacobian $|\Delta(t)|^2 =
\prod_{\alpha > 0} |2 \sin(\langle\alpha, t\rangle/2)|^2$ appears as a
multiplicative factor on the Cartan torus $T \subset G$. The Weyl
dimension formula states
\begin{equation}\label{eq:weyl_dim_universal}
   \dim V_\pi
   \;=\; \prod_{\alpha > 0}
   \frac{\langle \lambda_\pi + \rho,\, \alpha\rangle}
        {\langle\rho,\, \alpha\rangle},
\end{equation}
which has the same Vandermonde-style $\prod_{\alpha > 0}$ product
structure as $|\Delta(t)|^2$. In the dual-Coxeter normalisation,
these two products \emph{exactly cancel at leading order}, eliminating
all rank-, $|W|$-, and $N_+$-dependence (where $N_+$ is the number of
positive roots).

What is left is the pure $1$D Stein--Weiss constant $4/\pi$ from each
positive-root direction's Euler--Catalan asymptote
\begin{equation*}
   \sum_{d \text{ odd}} \frac{1}{d^2} \;=\; \frac{\pi^2}{8},
\end{equation*}
which is the same identity that gives the $4/\pi$ constant in
Paper~38's rank-$1$ Stein--Weiss calculation for $\SU(2)$
\cite[App.~A]{paper38}. The universal real-line Stein--Weiss
asymptotic
\begin{equation*}
   \int_{\T^1} F_n(\theta)\,|\theta|\,d\theta
   \;\sim\; \frac{4}{\pi} \cdot \frac{\log n}{n}
\end{equation*}
is then applied to each of the $r$ Cartan-torus directions, with the
Plancherel $\times$ Vandermonde cancellation ensuring that the
proportionality constant aggregates to the universal $4/\pi$ rather
than to a rank-dependent product.

\paragraph{Structural reading.} The central spectral Fej\'er kernel
sees its underlying real-line Stein--Weiss analysis through
Peter--Weyl, \emph{not} through the global geometry of $G/T$.
Whichever compact Lie group you build the kernel on, the Plancherel
weight $\sqrt{\dim V_\pi}$ and the Vandermonde Jacobian
$|\Delta(t)|^2$ both have $\prod_{\alpha > 0}$ structure tied to the
same root system, and they cancel against each other at leading order
in the dual-Coxeter normalisation. The Weyl group order $|W|$ and the
rank $r$ contribute factors to both products, and the factors cancel
against each other in the ratio
$\dim V_\pi^2 / |\Delta(t)|^2$ that controls the dominant
contribution to $\gamma_\Lammax$.

This cancellation is the natural structural origin of Reading~A
(rank-invariant $c(G) = 4/\pi$). Reading~B (Weyl-formula
$c(G) = 2|W(G)|/\pi^r$) would require the Weyl-group order to imprint
multiplicatively on the leading-order rate, which the cancellation
theorem forbids and which the empirical extraction at $\Sp(2)$ and
$G_2$ falsifies independently with margin $33\%$ and $52\%$
respectively (Theorem~\ref{thm:universal_constant}).

\subsection{Identification with the master Mellin engine}\label{sec:master_mellin}

The universal constant $4/\pi$ identifies with one specific sub-mechanism
of a master Mellin engine on the spectral triple, in the following
sense. The spectral-action heat kernel admits a Mellin transform
\begin{equation}\label{eq:master_mellin}
   \mathcal M[\Tr(\Dop^k e^{-t\Dop^2})](s)
   \;=\; \int_0^\infty t^{s-1} \Tr(\Dop^k e^{-t\Dop^2}) dt,
\end{equation}
where $k \in \{0, 1, 2\}$ indexes three structurally distinct
sub-mechanisms by which transcendentals (specifically powers of $\pi$
and Catalan-type Dirichlet $L$-values) enter spectral observables on
the triple. The $k = 0$ sub-mechanism corresponds to the
heat-kernel-as-Hopf-base-measure direction:\ here the leading
transcendental signature on any compact homogeneous space is the
volume factor of the natural Hopf-type fibration, which for $G$
itself is $\Vol(G/T) / \Vol(G)$ up to standard Haar normalisation
factors. In the dual-Coxeter normalisation this factor is, modulo a
universal Cesaro doubling, the constant $4/\pi$.

The universality of $4/\pi$ in Theorem~\ref{thm:main}(ii) is therefore
the spectral-triple-rate-controlled manifestation of the $k = 0$
sub-mechanism of the master Mellin engine, sharpened from a single-group
observation (Paper~38) to a class-wide theorem. The
GeoVac-internal scoping memo~\cite{unified_gh_scoping_memo} originally
proposed this identification as an open question; the present
theorem closes it positively across the rank-$2$ classical and
exceptional cases, with the rank-invariance of the leading constant
being the load-bearing experimental input.

% =====================================================================
\section{Discussion and open problems}\label{sec:discussion}
% =====================================================================

\subsection{Quantitative tightness and the log-power test}\label{sec:tightness}

Theorem~\ref{thm:main} gives the asymptotic constant $4/\pi$ in
$\gamma_\Lammax(G) \sim (4/\pi)\log\Lammax/\Lammax$ but not the
next-order constant. Numerical fits at the panel size used in this
work ($\Lammax^2 \in [1, 1000]$) display Stein--Weiss bias of $3$--$6\%$
typical and $10$--$25\%$ at smaller-asymptotic panels; the
verdict in favour of Reading~A vs Reading~B is robust at this scale
because the discriminator gap is $27$--$91\%$ across rank-$2$ groups,
many multiples of the bias.

A separate open question is the \emph{log-power}:\ at rank $r \ge 2$,
is the rate single-log $(\log\Lammax/\Lammax)$ or could it admit a
double-log $(\log^2\Lammax/\Lammax)$ refinement? This question is
\emph{numerically inconclusive} at the $\Lammax$ ranges accessible to
the verification panel:\ a $\log^r\Lammax/\Lammax$ ansatz with $r = 2$
fits the panel within the Stein--Weiss bias on $\SU(3)$, $\Sp(2)$, and
$G_2$, and a clean discriminator would require $\Lammax^2 \gtrsim 10^5$
which is computationally expensive. We have stated
Theorem~\ref{thm:main} in the single-log form as the simpler reading
consistent with the data; the analytic Stein--Weiss derivation at rank
$2$ would close this in a $1$--$2$ week sub-sprint. The structural
implications differ:\ if the rate is single-log universally, the M1
master-Mellin reading is sharpest; if it splits into a $\log^r$ family
at rank $r$, each log-power level would correspond to a separate
Cesaro-moment sub-tier.

\subsection{Rank-$3$ numerical extraction and compute scaling}
\label{sec:rank3_compute}

The Plancherel $\times$ Vandermonde cancellation theorem
(Theorem~\ref{thm:universal_constant_analytical}) proves $c(G) =
4/\pi$ rigorously at all ranks. Empirically, we have verified this at
$\SU(2)$ (analytical, Paper~38), $\SU(3)$ ($c = 1.243$, $2.4\%$
$A$-error), $\Sp(2)$ ($c = 1.087$, $14.6\%$ $A$-error), and $G_2$
($c = 1.177$, $7.6\%$ $A$-error). At $\SU(4)$ (rank $3$), the
numerical extraction is substantially noisier --- $c = 0.900$ at
$29.3\%$ $A$-error, with the compute budget terminating the panel at
$\Lammax^2 = 300$.

Two factors degrade the extraction at rank $3$. First, the rank-$3$
Weyl integration is genuinely $3$-dimensional, increasing the
per-point quadrature cost by roughly an order of magnitude over
rank~$2$;\ Cartan-torus rows at $\Lammax^2 \approx 400$ take
$15$--$30$ minutes each. Second, the pre-asymptotic structure is
stronger at rank~$3$ than at rank~$2$, making the leading-order
extraction more sensitive to the fit form. A forced-coefficient
$\chi^2$ analysis using a four-parameter subleading-aware ansatz
still favours Reading~A (consistency with $c = 4/\pi$) over Reading~B
(consistency with the Weyl-formula candidate $c(\SU(4)) =
48/\pi^3 \approx 1.548$) by a small margin ($4.6\%$ in
r.s.s.). The $\SU(4)$ extraction is therefore consistent with the
analytical universality theorem but does not independently confirm it
as cleanly as $\Sp(2)$ and $G_2$ at rank~$2$ did.

The decisive empirical falsification of Reading~B remains the $91\%$
$A$-vs-$B$ gap at $G_2$. The $\SU(4)$ data adds a third rank,
contributing to the cross-group consistency check, but does not
contribute decisively to the discriminator on its own. Detailed
extraction analysis is in the supporting memo~\cite{su4_rate_memo}.

\subsection{Extension to compact symmetric spaces $G/K$}\label{sec:symmetric}

The natural next-step generalisation is from compact Lie groups
(Class~1 in the scoping memo) to compact Riemannian symmetric spaces
$G/K$ (Class~2). The propinquity convergence theorem at full Class~2
generality would subsume Theorem~\ref{thm:main} as the ``group case''
$G \times G / G_{\mathrm{diag}}$ and would additionally cover spheres
$S^n = \SO(n+1)/\SO(n)$ for $n \ge 2$, complex projective spaces
$\mathbb{CP}^n = \SU(n+1)/(U(n) \times U(1))$, complex
Grassmannians, and full flag varieties.

The technical overhead is the $K$-invariance bookkeeping in the
Peter--Weyl decomposition $C^\infty(G/K) = \bigoplus_\pi V_\pi^K$ and
the corresponding zonal-spherical-function machinery. The central
spectral Fej\'er kernel construction extends to a $K$-bi-invariant
class-function kernel on $G$, and the Lipschitz comparison via the
Dirac-triangle inequality lifts to the $G/K$ setting via the
spherical functions of Helgason~\cite{helgason1984}. We expect the
universal rate constant $4/\pi$ to persist at the same level of
generality, but the verification at rank $\ge 2$ symmetric spaces is
deferred to a follow-on paper.

\subsection{Relation to Rieffel's coadjoint-orbit programme}\label{sec:rieffel}

A complementary direction in the quantum metric convergence literature
is Rieffel's coadjoint-orbit programme~\cite{rieffel2009, rieffel2023},
which constructs $G$-invariant Dirac operators on coadjoint orbits
$\mathcal{O} \subset \mathfrak{g}^*$ of a compact Lie group $G$ and
proves that fuzzy approximations (matrix algebras of growing
dimension) converge to $C(\mathcal{O})$ in spectral propinquity. This
work covers compact symmetric spaces that arise as coadjoint orbits
(notably $S^2 = \SU(2)/U(1)$, $\mathbb{CP}^n$, and complex Grassmannians)
and converges \emph{from below} (matrix algebras of growing dimension).

Theorem~\ref{thm:main} converges \emph{from above} (spectral truncation
of the continuum object) and treats compact Lie groups $G$ themselves
(which are not coadjoint orbits but rather principal homogeneous
spaces for $G$ acting on itself by left translation). The two
directions are categorically complementary:\ Rieffel's framework does
not cover $S^3 = \SU(2)$, and Theorem~\ref{thm:main} does not cover
$S^2 = \SU(2)/U(1)$. The natural unified statement, which would
subsume both, lives at the level of compact symmetric spaces $G/K$
discussed in \S\ref{sec:symmetric}.

\subsection{Lorentzian / Wick-rotated analog}\label{sec:lorentzian}

The Kostant cubic Dirac spectral triple is Riemannian, with
KO-dimension $\dim G \pmod 8$. The Lorentzian analog, via Wick rotation
of $G$ to a non-compact form, carries different KO-dimension and sign
convention, and may admit a different propinquity theory entirely. The
Bizi--Brouder--Besnard framework~\cite{bizi_brouder_besnard2018} for
$(m, n) \in (\Z/8)^2$ Krein-class spectral triples provides the
algebraic backbone for the Lorentzian extension, but the corresponding
quantum-Gromov--Hausdorff theory does not yet exist in the published
literature; a Lorentzian Latr\'emoli\`ere-style propinquity for
Krein-class triples remains open.

\subsection{Cross-manifold tensor products and the Coulomb/HO
asymmetry}\label{sec:cross_manifold}

A natural extension of Theorem~\ref{thm:main} would aim at tensor
products of spectral triples on different manifolds, in particular
$\Tcal_G \otimes \Tcal_H$ for compact Lie groups $G, H$ of possibly
different rank. The same-manifold tensor product
$\Tcal_{\SU(2)}^{\lambda_a} \otimes \Tcal_{\SU(2)}^{\lambda_b}$ at
distinct focal lengths is the subject of Paper~39~\cite{paper39}; the
genuine cross-manifold case $\Tcal_G \otimes \Tcal_H$ with $G \ne H$
should extend mechanically from Paper~39 with the Pythagorean Leibniz
constant generalising to a Cartan-product-type bound. We have not
written out the details.

\subsection{Connection to Hekkelman--McDonald non-commutative integral}
\label{sec:hekkelman_mcdonald}

Hekkelman and McDonald~\cite{hekkelman_mcdonald2024b} construct a
non-commutative integral on spectrally truncated triples and link it
to a Szeg\"o limit formula. Their construction lives on $T^d$ (abelian),
but the underlying truncation framework is the
Connes--van~Suijlekom paradigm common with the present paper. We expect
that the non-commutative integral of \cite{hekkelman_mcdonald2024b}
admits a class-wide compact-Lie-group extension under the same
Peter--Weyl machinery used here, with the rate constant $4/\pi$
appearing as the leading Szeg\"o-type asymptotic. The detailed
generalisation is open and would constitute a Szeg\"o companion to
the present propinquity theorem.

% =====================================================================
\section*{Acknowledgments}
% =====================================================================

The author thanks Alain Connes, Walter van~Suijlekom, Fran\c{c}ois
Latr\'emoli\`ere, Matilde Marcolli, Edward Hekkelman, Edward McDonald,
Malte Leimbach, Yvann Gaudillot-Estrada, and Marc Rieffel for the
foundational frameworks on which this paper builds; remaining errors
are the author's own. This work was developed in an AI-augmented
agentic workflow with Anthropic's Claude:\ implementation of the
internal proof-assistant verification machinery, drafting of the
internal proof memos used as supporting documentation, and
bibliographic collation were performed collaboratively. All theorem
statements, design choices, and editorial decisions are the author's
responsibility. The author thanks the open-source compact-Lie-group
representation-theory community (notably the sympy, SageMath, and
LiE teams) for the computational infrastructure used in the supporting
numerical verifications at $\SU(3), \SU(4), \Sp(2), G_2$.

% =====================================================================
\bibliographystyle{plainnat}
\begin{thebibliography}{99}

\bibitem[Agricola(2003)]{agricola2003}
I.~Agricola,
\newblock \emph{Connections on naturally reductive spaces, their Dirac
operator and homogeneous models in string theory},
\newblock Comm.~Math.~Phys. \textbf{232} (2003), 535--563.

\bibitem[Bizi, Brouder \& Besnard(2018)]{bizi_brouder_besnard2018}
N.~Bizi, C.~Brouder, and F.~Besnard,
\newblock \emph{Doubled spectral triples for Lorentzian signatures},
\newblock J.~Math.~Phys. \textbf{59} (2018), 062303;
\newblock arXiv:1611.07062.

\bibitem[Bo{\.z}ejko \& Fendler(1991)]{bozejko_fendler1991}
M.~Bo{\.z}ejko and G.~Fendler,
\newblock \emph{Herz--Schur multipliers and uniformly bounded
representations of discrete groups},
\newblock Arch.~Math. \textbf{57} (1991), 290--298.

\bibitem[Bourbaki(1975)]{bourbaki_lie_8}
N.~Bourbaki,
\newblock \emph{Groupes et alg\`ebres de Lie, Chapitres 7--9},
\newblock El\'ements de math\'ematique, Hermann/Springer, 1975
(English translation:\ Springer, 2005).

\bibitem[Branson(1996)]{branson1996}
T.~P.~Branson,
\newblock \emph{Stein--Weiss operators and ellipticity},
\newblock J.~Funct.~Anal. \textbf{151} (1997), 334--383.

\bibitem[Connes(1995)]{connes1995}
A.~Connes,
\newblock \emph{Noncommutative Geometry},
\newblock Academic Press, 1995.

\bibitem[Connes \& van~Suijlekom(2021)]{connes_vs2021}
A.~Connes and W.~D.~van~Suijlekom,
\newblock \emph{Spectral truncations in noncommutative geometry and
operator systems},
\newblock Comm.~Math.~Phys. \textbf{383} (2021), 2021--2067;
\newblock arXiv:2004.14115.

\bibitem[Fulton \& Harris(1991)]{fulton_harris1991}
W.~Fulton and J.~Harris,
\newblock \emph{Representation Theory:\ A First Course},
\newblock Graduate Texts in Mathematics \textbf{129},
Springer, 1991.

\bibitem[Gaudillot-Estrada \& van~Suijlekom(2025)]{gaudillot_estrada_vs2025}
Y.~Gaudillot-Estrada and W.~D.~van~Suijlekom,
\newblock \emph{Convergence of spectral truncations for compact metric
groups},
\newblock Int.~Math.~Res.~Not. (IMRN) \textbf{2025} (2025), rnaf197;
\newblock arXiv:2310.14733.

\bibitem[Hawkins(2000)]{hawkins2000}
E.~Hawkins,
\newblock \emph{Geometric quantization of vector bundles and the
correspondence with deformation quantization},
\newblock Comm.~Math.~Phys. \textbf{215} (2000), 409--432.

\bibitem[Hekkelman(2022)]{hekkelman2022}
E.~Hekkelman,
\newblock \emph{Truncated geometry on the circle},
\newblock Lett.~Math.~Phys. \textbf{112} (2022), 20;
\newblock arXiv:2111.13865.

\bibitem[Hekkelman \& McDonald(2024a)]{hekkelman_mcdonald2024}
E.~Hekkelman and E.~McDonald,
\newblock \emph{Spectral truncations of $T^d$ and quantum metric
geometry},
\newblock J.~Noncommut.~Geom., to appear (2024);
\newblock arXiv:2403.18619.

\bibitem[Hekkelman \& McDonald(2024b)]{hekkelman_mcdonald2024b}
E.~Hekkelman and E.~McDonald,
\newblock \emph{A noncommutative integral on spectrally truncated
spectral triples, and a link with quantum ergodicity},
\newblock J.~Funct.~Anal., to appear (2025);
\newblock arXiv:2412.00628.

\bibitem[Hekkelman, McDonald \& van~Suijlekom(2024)]{hekkelman_mcdonald_vs2024_ucp}
E.~Hekkelman, E.~McDonald, and W.~D.~van~Suijlekom,
\newblock \emph{UCP maps in spectral truncations of compact metric
spaces},
\newblock arXiv:2410.15454, 2024.

\bibitem[Helgason(1984)]{helgason1984}
S.~Helgason,
\newblock \emph{Groups and Geometric Analysis},
\newblock Academic Press, 1984.

\bibitem[Kac(1990)]{kac1990}
V.~G.~Kac,
\newblock \emph{Infinite-Dimensional Lie Algebras},
\newblock 3rd ed., Cambridge Univ.~Press, 1990.

\bibitem[Kostant(1999)]{kostant1999}
B.~Kostant,
\newblock \emph{A cubic Dirac operator and the emergence of Euler number
multiplets of representations for equal rank subgroups},
\newblock Duke Math.~J. \textbf{100} (1999), 447--501.

\bibitem[Latr\'emoli\`ere(2016)]{latremoliere2016}
F.~Latr\'emoli\`ere,
\newblock \emph{The dual Gromov--Hausdorff propinquity},
\newblock Banach J.~Math.~Anal. \textbf{10} (2016), 175--229.

\bibitem[Latr\'emoli\`ere(2017)]{latremoliere_metric_st_2017}
F.~Latr\'emoli\`ere,
\newblock \emph{The Gromov--Hausdorff propinquity for metric spectral
triples},
\newblock Adv.~Math. \textbf{404} (2022), Paper No.~108393, 56pp;
\newblock preprint arXiv:1811.10843, 2017.

\bibitem[Latr\'emoli\`ere(2018)]{latremoliere2018}
F.~Latr\'emoli\`ere,
\newblock \emph{The Gromov--Hausdorff propinquity},
\newblock Trans.~Amer.~Math.~Soc. \textbf{370} (2018), 365--411.

\bibitem[Leimbach(2024)]{leimbach2024}
M.~Leimbach,
\newblock \emph{Convergence of Peter--Weyl truncations of compact quantum
groups},
\newblock J.~Noncommut.~Geom., to appear (2024);
\newblock arXiv:2409.16698.

\bibitem[Leimbach \& van~Suijlekom(2024)]{leimbach_vs2024}
M.~Leimbach and W.~D.~van~Suijlekom,
\newblock \emph{Gromov--Hausdorff convergence of spectral truncations
of the torus},
\newblock Adv.~Math. \textbf{439} (2024), Paper No.~109496.

\bibitem[Rieffel(2009)]{rieffel2009}
M.~A.~Rieffel,
\newblock \emph{Dirac operators for coadjoint orbits of compact Lie
groups},
\newblock M\"unster J.~Math. \textbf{2} (2009), 265--298;
\newblock arXiv:0812.2884.

\bibitem[Rieffel(2023)]{rieffel2023}
M.~A.~Rieffel,
\newblock \emph{Dirac operators for matrix algebras converging to
coadjoint orbits},
\newblock Comm.~Math.~Phys. \textbf{401} (2023), 1951--2008;
\newblock arXiv:2108.01136.

\bibitem[Stein \& Weiss(1971)]{stein_weiss1971}
E.~M.~Stein and G.~Weiss,
\newblock \emph{Introduction to Fourier Analysis on Euclidean Spaces},
\newblock Princeton Math.~Ser. \textbf{32}, Princeton Univ.~Press, 1971.

\bibitem[Zygmund(2002)]{zygmund2002}
A.~Zygmund,
\newblock \emph{Trigonometric Series, Vol.~I and II},
\newblock 3rd ed., Cambridge Univ.~Press, 2002.

\bibitem[Loutey, Paper~38(2026)]{paper38}
J.~Loutey,
\newblock \emph{Latr\'emoli\`ere propinquity convergence of truncated
Camporesi--Higuchi spectral triples on $S^3$},
\newblock GeoVac internal preprint (May 2026), DOI via Zenodo.

\bibitem[Loutey, Paper~39(2026)]{paper39}
J.~Loutey,
\newblock \emph{Tensor-product propinquity convergence of metric
spectral triples on $S^3$ at distinct focal lengths},
\newblock GeoVac internal preprint (May 2026), DOI via Zenodo.

\bibitem[Loutey, internal memo (unified GH scoping)]{unified_gh_scoping_memo}
J.~Loutey,
\newblock \emph{Unified GH-convergence scoping memo}
(\texttt{debug/unified\_gh\_scoping\_memo.md}),
\newblock GeoVac internal working note (May 2026).

\bibitem[Loutey, internal memo (Sp(2) and G$_2$ rate constant)]{sp2_g2_rate_memo}
J.~Loutey,
\newblock \emph{$\Sp(2)$ and $G_2$ rate constant: A vs B discriminator test}
(\texttt{debug/sp2\_g2\_rate\_constant\_memo.md}),
\newblock GeoVac internal working note (May 2026).

\bibitem[Loutey, internal memo (Dirac-triangle)]{dirac_triangle_memo}
J.~Loutey,
\newblock \emph{Dirac-triangle inequality: full proof with PRV-summand
bound at all ranks}
(\texttt{debug/dirac\_triangle\_full\_proof.md}),
\newblock GeoVac internal working note (May 2026).

\bibitem[Loutey, internal memo (L2 universal rate)]{l2_universal_rate_memo}
J.~Loutey,
\newblock \emph{Analytical proof of the universal $4/\pi$ rate constant
via the Plancherel-weight $\times$ Vandermonde-Jacobian cancellation
theorem}
(\texttt{debug/l2\_universal\_rate\_proof.md}),
\newblock GeoVac internal working note (May 2026).

\bibitem[Loutey, internal memo (SU(4) rate constant)]{su4_rate_memo}
J.~Loutey,
\newblock \emph{SU(4) rate constant extraction at rank~$3$:
compute-budget analysis and forced-coefficient verdict}
(\texttt{debug/su4\_rate\_constant\_memo.md}),
\newblock GeoVac internal working note (May 2026).

\bibitem[Kumar(1988)]{kumar1988}
S.~Kumar,
\newblock Proof of the Parthasarathy--Ranga Rao--Varadarajan
conjecture,
\newblock \emph{Invent. Math.} \textbf{93} (1988), 117--130.

\bibitem[Mathieu(1989)]{mathieu1989}
O.~Mathieu,
\newblock Construction d'un groupe de Kac--Moody et applications,
\newblock \emph{Compos. Math.} \textbf{69} (1989), 37--60.

\bibitem[Polo(1994)]{polo1994}
P.~Polo,
\newblock Vari\'et\'es de Schubert et excellentes filtrations,
\newblock \emph{Ast\'erisque} \textbf{173--174} (1989), 281--311.

\bibitem[Vinberg(1990)]{vinberg1990}
\`E.~B.~Vinberg,
\newblock Some commutative subalgebras of a universal enveloping
algebra,
\newblock \emph{Math. USSR-Izv.} \textbf{36} (1991), 1--22 (translation
of \emph{Izv. Akad. Nauk SSSR Ser. Mat.} \textbf{54} (1990), 3--25).

\end{thebibliography}

\end{document}
